% =====================================================================
\chapter{Soluciones de los ejercicios}
\label{ap:soluciones}
% =====================================================================

Cada solución se presenta en cinco partes: datos, planteamiento, desarrollo, resultado y comprobación. En los ejercicios de demostración ($\star\star\star$), los pasos del desarrollo, leídos en orden, forman la demostración redactada tal como debería entregarse. La comprobación indica una forma independiente de verificar el resultado (un valor numérico, un caso particular, una gráfica) o señala dónde se ha usado cada hipótesis. Los valores numéricos se han comprobado con un ordenador y se redondean al número de cifras indicado.

% ---------------------------------------------------------------------
\section{Tema 1}

\begin{solucion}{eje:1-dominios}
\datos cuatro fórmulas.
\planteamiento dominio natural (\cref{def:funcion}): raíz cuadrada de argumento $\ge0$, logaritmo de argumento $>0$, denominador no nulo.
\pasos
\begin{enumerate}[label=(\alph*)]
  \item $x^{2}-5x+6=(x-2)(x-3)\ge0$ si y solo si $x\le2$ o $x\ge3$: $\dom f=(-\infty,2]\cup[3,\infty)$.
  \item Se necesita $\frac{1+x}{1-x}>0$ (y $x\ne1$): numerador y denominador del mismo signo. Ambos positivos: $-1<x<1$; ambos negativos: $x<-1$ y $x>1$, imposible. $\dom g=(-1,1)$.
  \item Se necesita $x>0$ y $\ln x\ne0$, es decir, $x\ne1$: $\dom h=(0,1)\cup(1,\infty)$.
  \item Se necesita $x>0$ y $\ln x\ge0$, es decir, $x\ge1$: $\dom k=[1,\infty)$.
\end{enumerate}
\resultado (a) $(-\infty,2]\cup[3,\infty)$; (b) $(-1,1)$; (c) $(0,1)\cup(1,\infty)$; (d) $[1,\infty)$.
\comprobacion valores de prueba: $f(2{,}5)=\sqrt{-0{,}25}$ no existe; $g(0)=\ln1=0$ existe; $h(1)=1/0$ no existe; $k(0{,}5)=\sqrt{-0{,}69}$ no existe.
\end{solucion}

\begin{solucion}{eje:1-transformaciones}
\datos $(0,0)$, $(1,2)$, $(3,-1)$ están en la gráfica de $f$.
\planteamiento si $(x_0,y_0)$ está en la gráfica de $f$, se busca el punto correspondiente en cada transformación (\cref{prop:transformaciones}).
\pasos
\begin{enumerate}[label=(\alph*)]
  \item $y=f(x-2)+1$: traslación de 2 a la derecha y 1 hacia arriba; $(x_0,y_0)\mapsto(x_0+2,y_0+1)$: $(2,1)$, $(3,3)$, $(5,0)$.
  \item $y=3f(x)$: dilatación vertical de factor~3; $(x_0,3y_0)$: $(0,0)$, $(1,6)$, $(3,-3)$.
  \item $y=f(2x)$: compresión horizontal de factor~2; $(x_0/2,y_0)$: $(0,0)$, $(\tfrac12,2)$, $(\tfrac32,-1)$.
  \item $y=-f(-x)$: simetría respecto del origen; $(-x_0,-y_0)$: $(0,0)$, $(-1,-2)$, $(-3,1)$.
\end{enumerate}
\resultado los puntos indicados en cada apartado.
\comprobacion en (c), $x=\tfrac12$ da $y=f(2\cdot\tfrac12)=f(1)=2$; en (d), $x=-1$ da $y=-f(1)=-2$.
\end{solucion}

\begin{solucion}{eje:1-paridad}
\datos seis funciones.
\planteamiento calcular $f(-x)$ y compararlo con $f(x)$ y $-f(x)$ (\cref{def:propiedades}); para descartar ambas propiedades basta un valor.
\pasos
\begin{enumerate}[label=(\alph*)]
  \item $(-x)^{3}-(-x)=-(x^{3}-x)$: \textbf{impar}.
  \item $\abs{-x}+\cos(-x)=\abs{x}+\cos x$: \textbf{par}.
  \item $e^{-1}\ne e^{1}$ y $e^{-1}\ne-e^{1}$: \textbf{ninguna}.
  \item $\ln\frac{1+x}{1-x}=-\ln\frac{1-x}{1+x}$: \textbf{impar}.
  \item $(-x)\sigma(-x)=-x(1-\sigma(x))=x\sigma(x)-x$, que no es $\pm x\sigma(x)$; con $x=1$: $f(1)=\sigma(1)\approx0{,}731$ y $f(-1)=-\sigma(-1)\approx-0{,}269$. \textbf{Ninguna}.
  \item Por la \cref{prop:relu}(ii), $\softplus(-x)=\softplus(x)-x$, luego $\softplus(-x)-\frac{-x}{2}=\softplus(x)-x+\frac x2=\softplus(x)-\frac x2$: \textbf{par}.
\end{enumerate}
\resultado impares: (a), (d); pares: (b), (f); ninguna: (c), (e).
\comprobacion (f) en $x=\pm2$: $\softplus(2)-1=2{,}1269-1=1{,}1269$ y $\softplus(-2)+1=0{,}1269+1=1{,}1269$.
\end{solucion}

\begin{solucion}{eje:1-composicion}
\datos $f(x)=\sqrt x$ ($\dom f=[0,\infty)$), $g(x)=1-x^{2}$ ($\dom g=\R$).
\planteamiento $\dom(u\circ v)=\{x\in\dom v: v(x)\in\dom u\}$.
\pasos
\begin{enumerate}[label=\arabic*.]
  \item $(f\circ g)(x)=\sqrt{1-x^{2}}$; se necesita $1-x^{2}\ge0$: dominio $[-1,1]$.
  \item $(g\circ f)(x)=1-(\sqrt x)^{2}=1-x$, pero solo para $x\in\dom f=[0,\infty)$.
  \item $(f\circ f)(x)=\sqrt{\sqrt x}=x^{1/4}$, dominio $[0,\infty)$.
  \item $(g\circ g)(x)=1-(1-x^{2})^{2}=2x^{2}-x^{4}$, dominio $\R$.
\end{enumerate}
\resultado $f\circ g=\sqrt{1-x^{2}}$ en $[-1,1]$; $g\circ f=1-x$ en $[0,\infty)$; $f\circ f=x^{1/4}$ en $[0,\infty)$; $g\circ g=2x^{2}-x^{4}$ en $\R$.
\comprobacion $g\circ f$ no es la función $1-x$ en $\R$: en $x=-1$ la composición no está definida, aunque la fórmula simplificada sí lo esté. Valor de prueba: $(g\circ g)(1)=g(0)=1=2-1$.
\end{solucion}

\begin{solucion}{eje:1-ecuaciones}
\datos cuatro ecuaciones.
\planteamiento aislar la exponencial y tomar logaritmos; en las logarítmicas, usar $\ln a+\ln b=\ln(ab)$ y comprobar el dominio.
\pasos
\begin{enumerate}[label=(\alph*)]
  \item $e^{3x}=15$, $3x=\ln15$, $x=\frac{\ln15}{3}\approx0{,}9027$.
  \item $(3x-2)\ln7=\ln11$, $x=\frac13\bigl(2+\frac{\ln11}{\ln7}\bigr)\approx1{,}0774$.
  \item Se necesita $x>2$. $\ln(x(x-2))=\ln4$, y como $\ln$ es inyectiva, $x^{2}-2x-4=0$, $x=1\pm\sqrt5$. Solo $1+\sqrt5\approx3{,}236$ cumple $x>2$.
  \item $4^{x+1}=32$, $2^{2x+2}=2^{5}$, $2x+2=5$, $x=\tfrac32$.
\end{enumerate}
\resultado (a) $\frac{\ln15}{3}$; (b) $\frac{2}{3}+\frac{\ln11}{3\ln7}$; (c) $1+\sqrt5$; (d) $\tfrac32$.
\comprobacion (c): $(1+\sqrt5)(\sqrt5-1)=5-1=4$. (d): $4^{2{,}5}=2^{5}=32$.
\end{solucion}

\begin{solucion}{eje:1-hiperbolicas-calculo}
\datos definiciones de $\sinh$ y $\cosh$.
\planteamiento sustituir la \cref{def:hiperbolicas} y usar la \cref{prop:hiperbolicas}(iii).
\pasos
\begin{enumerate}[label=(\alph*)]
  \item $2\cosh(\ln x)=e^{\ln x}+e^{-\ln x}=x+\frac1x$.
  \item $\cosh4x+\sinh4x=e^{4x}$.
  \item $\cosh2x-\sinh2x=e^{-2x}$.
  \item $\ln(e^{x})+\ln(e^{-x})=x-x=0$ (o bien $\ln(\cosh^{2}x-\sinh^{2}x)=\ln1=0$).
\end{enumerate}
\resultado (a) $x+1/x$; (b) $e^{4x}$; (c) $e^{-2x}$; (d) $0$.
\comprobacion (a) con $x=2$: $2\cosh(\ln2)=2\cdot1{,}25=2{,}5=2+\tfrac12$.
\end{solucion}

\begin{solucion}{eje:1-sigmoide-valores}
\datos $\sigma(x)=1/(1+e^{-x})$, $\logit(p)=\ln\frac{p}{1-p}$.
\planteamiento evaluar directamente y usar la inversa (\cref{prop:sigmoide}(iii)).
\pasos
\begin{enumerate}[label=(\alph*)]
  \item $\sigma(0)=\frac12$; $\sigma(\ln3)=\frac{1}{1+1/3}=\frac34$; $\sigma(-\ln3)=\frac{1}{1+3}=\frac14$.
  \item $x=\logit(0{,}9)=\ln9\approx2{,}197$; $x=\logit(0{,}01)=\ln\frac{1}{99}=-\ln99\approx-4{,}595$.
  \item $\sigma$ es estrictamente creciente, luego $\sigma(x)>0{,}99\iff x>\logit(0{,}99)=\ln99\approx4{,}595$.
  \item $\tanh(\ln2)=\frac{2-1/2}{2+1/2}=\frac{3/2}{5/2}=\frac35$.
\end{enumerate}
\resultado (a) $\tfrac12$, $\tfrac34$, $\tfrac14$; (b) $\ln9$ y $-\ln99$; (c) $x>\ln99$; (d) $\tfrac35$.
\comprobacion $\sigma(\ln3)+\sigma(-\ln3)=1$, como exige la simetría $\sigma(-x)=1-\sigma(x)$. Con (iv) de la \cref{prop:sigmoide}: $2\sigma(2\ln2)-1=2\sigma(\ln4)-1=2\cdot\frac45-1=\frac35$.
\end{solucion}

\begin{solucion}{eje:1-inversa-racional}
\datos $f(x)=\frac{x+1}{x-1}$, $x\ne1$.
\planteamiento componer y simplificar; una función con $f\circ f=\mathrm{id}$ es biyectiva y su propia inversa.
\pasos
\begin{enumerate}[label=\arabic*.]
  \item Para $x\ne1$, $f(x)\ne1$ (si $\frac{x+1}{x-1}=1$, entonces $x+1=x-1$, absurdo), luego $f\circ f$ está definida en $\R\setminus\{1\}$.
  \item $f(f(x))=\dfrac{\frac{x+1}{x-1}+1}{\frac{x+1}{x-1}-1}=\dfrac{(x+1)+(x-1)}{(x+1)-(x-1)}=\dfrac{2x}{2}=x$.
  \item Como $f\circ f=\mathrm{id}$, $f$ es inyectiva (si $f(x)=f(y)$, aplicando $f$, $x=y$) y sobreyectiva sobre $\R\setminus\{1\}$ (todo $y\ne1$ es $f(f(y))$). Luego $f^{-1}=f$ y la imagen es $\R\setminus\{1\}$.
  \item La gráfica de $f^{-1}$ es la simétrica de la de $f$ respecto de $y=x$ (\cref{prop:inversa-monotona}; el argumento de la simetría no usa la monotonía); como $f^{-1}=f$, la gráfica es simétrica de sí misma.
\end{enumerate}
\resultado $f\circ f=\mathrm{id}$, $f^{-1}=f$, imagen $\R\setminus\{1\}$; gráfica simétrica respecto de $y=x$.
\comprobacion $f(3)=2$ y $f(2)=3$: el punto $(3,2)$ y su simétrico $(2,3)$ están en la gráfica.
\end{solucion}

\begin{solucion}{eje:1-sigmoide-simetria}
\datos $\sigma(x)=\frac{1}{1+e^{-x}}$, $\tanh x=\frac{e^{x}-e^{-x}}{e^{x}+e^{-x}}$.
\planteamiento cálculo algebraico directo con exponenciales.
\pasos
\begin{enumerate}[label=\arabic*.]
  \item Sea $x\in\R$. Multiplicando numerador y denominador por $e^{-x}$, $\sigma(-x)=\frac{1}{1+e^{x}}=\frac{e^{-x}}{e^{-x}+1}$.
  \item Por otra parte, $1-\sigma(x)=\frac{1+e^{-x}-1}{1+e^{-x}}=\frac{e^{-x}}{1+e^{-x}}$, que coincide con lo anterior. Luego $\sigma(-x)=1-\sigma(x)$.
  \item Para (iv): $2\sigma(2x)-1=\frac{2}{1+e^{-2x}}-1=\frac{1-e^{-2x}}{1+e^{-2x}}$. Multiplicando numerador y denominador por $e^{x}$ se obtiene $\frac{e^{x}-e^{-x}}{e^{x}+e^{-x}}=\tanh x$.
  \item Sustituyendo $x$ por $x/2$ en la identidad anterior, $\tanh\frac x2=2\sigma(x)-1$, y despejando, $\sigma(x)=\frac12\bigl(1+\tanh\frac x2\bigr)$. \hfill$\blacksquare$
\end{enumerate}
\resultado $\sigma(-x)=1-\sigma(x)$, $\tanh x=2\sigma(2x)-1$ y $\sigma(x)=\frac12(1+\tanh\frac x2)$.
\comprobacion con $x=\ln3$: $\frac12(1+\tanh\frac{\ln3}{2})=\frac12\bigl(1+\frac{\sqrt3-1/\sqrt3}{\sqrt3+1/\sqrt3}\bigr)=\frac12(1+\frac{2}{4})=\frac34=\sigma(\ln3)$.
\end{solucion}

\begin{solucion}{eje:1-hiperbolicas-suma}
\datos $\cosh u\pm\sinh u=e^{\pm u}$ (\cref{prop:hiperbolicas}(iii)).
\planteamiento desarrollar el miembro derecho con exponenciales; para (c), dividir (a) entre (b).
\pasos
\begin{enumerate}[label=\arabic*.]
  \item Sean $x,y\in\R$. $\sinh x\cosh y+\cosh x\sinh y=\frac14\bigl[(e^{x}-e^{-x})(e^{y}+e^{-y})+(e^{x}+e^{-x})(e^{y}-e^{-y})\bigr]$.
  \item Desarrollando: el primer producto es $e^{x+y}+e^{x-y}-e^{-x+y}-e^{-x-y}$ y el segundo $e^{x+y}-e^{x-y}+e^{-x+y}-e^{-x-y}$. La suma es $2e^{x+y}-2e^{-(x+y)}$, y al dividir por~4 se obtiene $\sinh(x+y)$. Esto prueba (a).
  \item Análogamente, $\cosh x\cosh y+\sinh x\sinh y=\frac14\bigl[(e^{x+y}+e^{x-y}+e^{-x+y}+e^{-x-y})+(e^{x+y}-e^{x-y}-e^{-x+y}+e^{-x-y})\bigr]=\frac{e^{x+y}+e^{-(x+y)}}{2}=\cosh(x+y)$. Esto prueba (b).
  \item Por (a) y (b), $\tanh(x+y)=\frac{\sinh x\cosh y+\cosh x\sinh y}{\cosh x\cosh y+\sinh x\sinh y}$. Dividiendo numerador y denominador por $\cosh x\cosh y>0$ se obtiene $\frac{\tanh x+\tanh y}{1+\tanh x\tanh y}$, que prueba (c). El denominador no se anula porque $\abs{\tanh x\tanh y}<1$. \hfill$\blacksquare$
\end{enumerate}
\resultado las tres identidades. Frente a las trigonométricas, $\sinh$ se comporta como $\sin$, pero en $\cosh(x+y)$ y en el denominador de $\tanh(x+y)$ aparece un signo $+$ donde las trigonométricas tienen $-$.
\comprobacion con $y=x$: (b) da $\cosh2x=\cosh^{2}x+\sinh^{2}x$; en $x=0$, $1=1+0$. (c) con $x=y=\ln2$: $\tanh(\ln4)=\frac{4-1/4}{4+1/4}=\frac{15}{17}$ y $\frac{2\cdot\frac35}{1+\frac{9}{25}}=\frac{6/5}{34/25}=\frac{15}{17}$.
\end{solucion}

\begin{solucion}{eje:1-arsinh}
\datos $\sinh x=\frac{e^{x}-e^{-x}}{2}$, $\tanh x=\frac{e^{2x}-1}{e^{2x}+1}$.
\planteamiento resolver $f(x)=y$ con el cambio $u=e^{x}>0$ (o $u=e^{2x}$); una solución única para cada $y$ prueba la biyectividad.
\pasos
\begin{enumerate}[label=\arabic*.]
  \item (a) Sea $y\in\R$. Con $u=e^{x}>0$, $\sinh x=y$ equivale a $u-\frac1u=2y$, es decir, $u^{2}-2yu-1=0$, cuyas raíces son $u=y\pm\sqrt{y^{2}+1}$.
  \item Como $\sqrt{y^{2}+1}>\abs y$, la raíz $y-\sqrt{y^{2}+1}$ es negativa y $y+\sqrt{y^{2}+1}$ es positiva. Hay, pues, exactamente un $u>0$ y exactamente un $x=\ln u$: $\sinh$ es biyectiva y $\arsinh y=\ln\bigl(y+\sqrt{y^{2}+1}\bigr)$.
  \item (b) Sea $y\in\R$. Con $u=e^{2x}>0$, $\tanh x=y$ equivale a $u-1=y(u+1)$, es decir, $u(1-y)=1+y$. Si $y=1$ no hay solución; si $y\ne1$, $u=\frac{1+y}{1-y}$, que es positivo si y solo si $-1<y<1$.
  \item Por tanto, cada $y\in(-1,1)$ tiene exactamente una preimagen $x=\frac12\ln\frac{1+y}{1-y}$, y ningún $y\notin(-1,1)$ la tiene: $\tanh\colon\R\to(-1,1)$ es biyectiva con la inversa del enunciado.
  \item (c) Con $p=\frac{1+y}{2}$, $1-p=\frac{1-y}{2}$ y $\frac{p}{1-p}=\frac{1+y}{1-y}$, luego $\frac12\logit(p)=\frac12\ln\frac{1+y}{1-y}=\artanh y$. \hfill$\blacksquare$
\end{enumerate}
\resultado $\arsinh y=\ln(y+\sqrt{y^{2}+1})$ en $\R$; $\artanh y=\frac12\ln\frac{1+y}{1-y}$ en $(-1,1)$; $\artanh y=\frac12\logit\frac{1+y}{2}$.
\comprobacion $\arsinh(\frac34)=\ln(\frac34+\frac54)=\ln2$ y $\sinh(\ln2)=\frac{2-1/2}{2}=\frac34$. $\artanh(\frac35)=\frac12\ln\frac{8/5}{2/5}=\frac12\ln4=\ln2$, coherente con el \cref{eje:1-sigmoide-valores}(d).
\end{solucion}

\begin{solucion}{eje:1-estricta}
\datos $f\colon[a,b]\to\R$ estrictamente creciente.
\planteamiento la inyectividad, por casos; la segunda afirmación, por contrarrecíproco.
\pasos
\begin{enumerate}[label=\arabic*.]
  \item Sean $x\ne y$ en $[a,b]$. Si $x<y$, $f(x)<f(y)$; si $y<x$, $f(y)<f(x)$. En ambos casos $f(x)\ne f(y)$: $f$ es inyectiva.
  \item Supóngase $f(x)<f(y)$ y, por reducción al absurdo, $x\ge y$. Si $x=y$, $f(x)=f(y)$; si $x>y$, $f(x)>f(y)$. En ambos casos $f(x)\ge f(y)$, contradicción. Luego $x<y$.
  \item Ejemplo: $g\colon[0,2]\to\R$ con $g(x)=x$ en $[0,1]$ y $g(x)=4-x$ en $(1,2]$. Cada tramo es inyectivo, y los valores del primero, $[0,1]$, y del segundo, $[2,3)$, son disjuntos; luego $g$ es inyectiva. Pero $0<1{,}5<2$ y $g(0)=0<g(1{,}5)=2{,}5>g(2)=2$, de modo que $g$ no es creciente ni decreciente.
\end{enumerate}
\resultado $f$ es inyectiva y $f(x)<f(y)\Rightarrow x<y$. La función $g(x)=x$ en $[0,1]$, $g(x)=4-x$ en $(1,2]$ es inyectiva y no monótona.
\comprobacion la $g$ del ejemplo es discontinua en $x=1$; el \cref{eje:3-inyectiva-monotona} prueba que una función \emph{continua} e inyectiva en un intervalo sí es estrictamente monótona, lo que explica que el contraejemplo necesite un salto.
\end{solucion}

\begin{solucion}{eje:1-monotonia-composicion}
\datos $f,g\colon\R\to\R$.
\planteamiento aplicar las definiciones a $x<y$ arbitrarios.
\pasos
\begin{enumerate}[label=\arabic*.]
  \item (a) Sean $x<y$. Como $f$ es creciente, $f(x)\le f(y)$; como $g$ es creciente, $g(f(x))\le g(f(y))$. Luego $g\circ f$ es creciente.
  \item (b) Sean $x<y$. Entonces $f(x)\le f(y)$ y, como $g$ es decreciente, $g(f(x))\ge g(f(y))$: $g\circ f$ es decreciente.
  \item (c) Si $x<y$, $f(x)\le f(y)$ y $g(x)\le g(y)$; sumando, $(f+g)(x)\le(f+g)(y)$.
  \item (d) Contraejemplo: $f(x)=g(x)=x$ son crecientes, pero $f(x)g(x)=x^{2}$ no lo es ($(-1)^{2}>0^{2}$).
  \item Aplicación: $x\mapsto x^{3}$ es creciente y $x\mapsto\sigma(-x)$ es decreciente (es $\sigma$, creciente, compuesta con $-x$, decreciente, que por un argumento como (b) da una función decreciente). Por (b), $x\mapsto\sigma(-x^{3})$ es decreciente. \hfill$\blacksquare$
\end{enumerate}
\resultado (a)--(c) demostradas; (d) falsa en general; $\sigma(-x^{3})$ es decreciente.
\comprobacion $\sigma(-0^{3})=0{,}5>\sigma(-1)=0{,}269$. En (d), el producto de crecientes \emph{positivas} sí es creciente: si $0<f(x)\le f(y)$ y $0<g(x)\le g(y)$, entonces $f(x)g(x)\le f(y)g(y)$.
\end{solucion}

\begin{solucion}{eje:1-par-impar}
\datos $D$ simétrico, $f\colon D\to\R$.
\planteamiento unicidad: suponer $f=p+q$ y obtener $p$ y $q$; existencia: comprobar que las fórmulas obtenidas funcionan.
\pasos
\begin{enumerate}[label=\arabic*.]
  \item \emph{Unicidad.} Supóngase $f=p+q$ con $p$ par y $q$ impar. Para $x\in D$, $f(x)=p(x)+q(x)$ y $f(-x)=p(x)-q(x)$. Sumando y restando, $p(x)=\frac{f(x)+f(-x)}{2}$ y $q(x)=\frac{f(x)-f(-x)}{2}$: $p$ y $q$ quedan determinadas por $f$.
  \item \emph{Existencia.} Se definen $p$ y $q$ por esas fórmulas. Entonces $p(-x)=\frac{f(-x)+f(x)}{2}=p(x)$, $q(-x)=\frac{f(-x)-f(x)}{2}=-q(x)$ y $p+q=f$.
  \item Para $f(x)=e^{x}$: $p=\frac{e^{x}+e^{-x}}{2}=\cosh x$ y $q=\sinh x$.
  \item Para $f=\softplus$: por la \cref{prop:relu}(ii), $f(x)-f(-x)=x$, luego $q(x)=\frac x2$, y $p(x)=\softplus(x)-\frac x2$, que es la función par del \cref{eje:1-paridad}(f). \hfill$\blacksquare$
\end{enumerate}
\resultado $p(x)=\frac{f(x)+f(-x)}{2}$, $q(x)=\frac{f(x)-f(-x)}{2}$; para $e^{x}$, $\cosh$ y $\sinh$; para la softplus, $\softplus(x)-\frac x2$ y $\frac x2$.
\comprobacion $\cosh1+\sinh1=1{,}5431+1{,}1752=2{,}7183=e$.
\end{solucion}

\begin{solucion}{eje:1-softplus}
\datos $\softplus(x)=\ln(1+e^{x})$, $\sigma(x)=1/(1+e^{-x})$.
\planteamiento manipular el argumento del logaritmo; para la monotonía, componer funciones crecientes.
\pasos
\begin{enumerate}[label=\arabic*.]
  \item $\softplus(x)=\ln\bigl(e^{x}(e^{-x}+1)\bigr)=x+\ln(1+e^{-x})=x+\softplus(-x)$.
  \item $-\ln\sigma(-x)=-\ln\frac{1}{1+e^{x}}=\ln(1+e^{x})=\softplus(x)$.
  \item Monotonía: $x\mapsto e^{x}$ es estrictamente creciente, luego también $1+e^{x}$, y $\ln$ es estrictamente creciente; la composición de estrictamente crecientes es estrictamente creciente (como en el \cref{eje:1-monotonia-composicion}(a), con desigualdades estrictas).
  \item Cambiando $x$ por $-x$ en el paso~2: $\softplus(-x)=-\ln\sigma(x)$, es decir, $\ln\sigma(x)=-\softplus(-x)$.
  \item Utilidad numérica: para $x=-1000$, $\sigma(x)=1/(1+e^{1000})$ es menor que el menor número representable en coma flotante y se redondea a~0, y $\ln0=-\infty$. En cambio, $-\softplus(1000)=-(1000+\softplus(-1000))=-(1000+\ln(1+e^{-1000}))$ se calcula sin problemas y vale $-1000$ con precisión de máquina. \hfill$\blacksquare$
\end{enumerate}
\resultado las identidades del enunciado; $\ln\sigma(x)=-\softplus(-x)$ evita el desbordamiento por defecto.
\comprobacion con $x=0$: $\ln\sigma(0)=\ln\frac12=-\ln2=-\softplus(0)$.
\end{solucion}
% ---------------------------------------------------------------------
\section{Tema 2}

\begin{solucion}{eje:2-negacion}
\datos las definiciones~\ref{def:limite-sucesion} y~\ref{def:limite-funcion}.
\planteamiento negar cuantificadores: $\neg\forall=\exists\neg$, $\neg\exists=\forall\neg$, y $\neg(P\Rightarrow Q)\equiv P\wedge\neg Q$ (módulo~0.1).
\pasos
\begin{enumerate}[label=\arabic*.]
  \item $a_n\not\to L$: $\exists\varepsilon>0\ \forall N\in\N\ \exists n\ge N:\ \abs{a_n-L}\ge\varepsilon$.
  \item $\lim_{x\to a}f(x)\ne L$: $\exists\varepsilon>0\ \forall\delta>0\ \exists x:\ 0<\abs{x-a}<\delta\ \wedge\ \abs{f(x)-L}\ge\varepsilon$.
  \item Para probar $a_n\not\to L$ hay que exhibir un $\varepsilon>0$ concreto y mostrar que, por grande que sea $N$, hay términos con índice $\ge N$ a distancia al menos $\varepsilon$ de $L$; equivalentemente, que hay infinitos términos fuera del intervalo $(L-\varepsilon,L+\varepsilon)$.
\end{enumerate}
\resultado las dos negaciones anteriores.
\comprobacion para $a_n=(-1)^{n}$ y $L=1$, $\varepsilon=1$ sirve: todos los términos de índice impar valen $-1$ y distan~2 de~1.
\end{solucion}

\begin{solucion}{eje:2-N}
\datos $a_n=\frac{2n}{n+3}$, $L=2$.
\planteamiento calcular $\abs{a_n-2}$ exactamente y despejar $n$.
\pasos
\begin{enumerate}[label=\arabic*.]
  \item $\abs{a_n-2}=\abs{\frac{2n-2n-6}{n+3}}=\frac{6}{n+3}$, que es decreciente en $n$.
  \item $\frac{6}{n+3}<0{,}01\iff n+3>600\iff n>597$. El menor $N$ es $598$.
  \item $\frac{6}{n+3}<10^{-6}\iff n+3>6\cdot10^{6}\iff n>5\,999\,997$. El menor $N$ es $5\,999\,998$.
\end{enumerate}
\resultado $N=598$ para $\varepsilon=0{,}01$ y $N=5\,999\,998$ para $\varepsilon=10^{-6}$.
\comprobacion $\abs{a_{597}-2}=6/600=0{,}01$, que no es $<0{,}01$; $\abs{a_{598}-2}=6/601<0{,}01$. Como la cota es decreciente, todos los $n\ge N$ sirven.
\end{solucion}

\begin{solucion}{eje:2-sucesiones-calculo}
\datos seis sucesiones.
\planteamiento dividir por la potencia dominante, multiplicar por el conjugado, usar el sándwich o la \cref{prop:notables-exp}(iii).
\pasos
\begin{enumerate}[label=(\alph*)]
  \item Dividiendo por $n^{2}$: $\frac{3-1/n+4/n^{2}}{5+2/n^{2}}\to\frac35$ (\cref{teo:algebra-sucesiones} y $1/n\to0$).
  \item $\sqrt{n^{2}+n}-n=\frac{(n^{2}+n)-n^{2}}{\sqrt{n^{2}+n}+n}=\frac{n}{\sqrt{n^{2}+n}+n}=\frac{1}{\sqrt{1+1/n}+1}\to\frac12$, usando la continuidad de la raíz (\cref{tema:3}).
  \item $\frac{2^{n}+3^{n}}{3^{n+1}}=\frac13\bigl(\frac23\bigr)^{n}+\frac13\to\frac13$ por la \cref{prop:geometrica}.
  \item $-\frac1n\le\frac{\sin n}{n}\le\frac1n$; por el sándwich, el límite es~0.
  \item Por la \cref{prop:notables-exp}(iii) con $t=-2$: $e^{-2}\approx0{,}1353$.
  \item $3^{n}\le2^{n}+3^{n}\le2\cdot3^{n}$, luego $3\le\sqrt[n]{2^{n}+3^{n}}\le3\cdot2^{1/n}$. Como $2^{1/n}\to1$ (\cref{eje:2-raizn}(b)), el sándwich da~3.
\end{enumerate}
\resultado (a) $\tfrac35$; (b) $\tfrac12$; (c) $\tfrac13$; (d) $0$; (e) $e^{-2}$; (f) $3$.
\comprobacion numérica: para $n=10^{6}$, (b) vale $0{,}49999988$ y (e) vale $0{,}1353350$, frente a $e^{-2}=0{,}1353353$.
\end{solucion}

\begin{solucion}{eje:2-funciones-calculo}
\datos cinco límites de funciones racionales o con raíces.
\planteamiento factorizar y simplificar el factor que se anula; en $\infty$, dividir por la potencia dominante; en las raíces, multiplicar por el conjugado.
\pasos
\begin{enumerate}[label=(\alph*)]
  \item $\frac{(x-2)(x+2)}{(x-1)(x-2)}=\frac{x+2}{x-1}$ para $x\ne2$; el límite es $\frac41=4$.
  \item $\frac{x+1}{(x-1)(x+1)}=\frac{1}{x-1}$ para $x\ne\pm1$. Cuando $x\to1^{+}$, $x-1\to0^{+}$ y el cociente tiende a $+\infty$; cuando $x\to1^{-}$, a $-\infty$. El límite (bilateral) no existe.
  \item $\frac{2-1/x^{2}}{1+4/x}\to2$.
  \item $\frac{\sqrt{1+x}-1}{x}=\frac{(1+x)-1}{x(\sqrt{1+x}+1)}=\frac{1}{\sqrt{1+x}+1}\to\frac12$.
  \item $x^{3}+8=(x+2)(x^{2}-2x+4)$, luego el cociente es $x^{2}-2x+4\to4+4+4=12$.
\end{enumerate}
\resultado (a) $4$; (b) $+\infty$ y $-\infty$; (c) $2$; (d) $\tfrac12$; (e) $12$.
\comprobacion numérica: en $x=2{,}000001$, (a) vale $3{,}999997$; en $x=10^{-6}$, (d) vale $0{,}49999988$.
\end{solucion}

\begin{solucion}{eje:2-notables}
\datos la \cref{tab:notables}.
\planteamiento reducir cada límite a uno conocido mediante un cambio de variable o un producto de factores con límite conocido; demostrar primero las entradas de la tabla que faltan.
\pasos
\begin{enumerate}[label=\arabic*.]
  \item \emph{Entradas pendientes.} $\frac{\tan x}{x}=\frac{\sin x}{x}\cdot\frac{1}{\cos x}\to1$. $\frac{a^{x}-1}{x}=\ln a\cdot\frac{e^{u}-1}{u}$ con $u=x\ln a\to0$ (si $a\ne1$), que tiende a $\ln a$ por~\eqref{eq:lim-exp}. Para $p\ne0$, $\frac{(1+x)^{p}-1}{x}=\frac{e^{u}-1}{u}\cdot\frac{p\ln(1+x)}{x}$ con $u=p\ln(1+x)\to0$, que tiende a $1\cdot p$ por la \cref{prop:notables-exp}(i). $\frac{\tanh x}{x}=\frac{e^{2x}-1}{2x}\cdot\frac{2}{e^{2x}+1}\to1\cdot\frac22=1$.
  \item (a) $\frac{\sin5x}{3x}=\frac53\cdot\frac{\sin5x}{5x}\to\frac53$.
  \item (b) $1$, por el paso~1.
  \item (c) $\frac{e^{2x}-1}{\sin x}=2\cdot\frac{e^{2x}-1}{2x}\cdot\frac{x}{\sin x}\to2$.
  \item (d) $\ln2$, por el paso~1 con $a=2$.
  \item (e) $1$, por el paso~1.
  \item (f) Con $t=1/x\to0^{+}$: $x\sin(1/x)=\frac{\sin t}{t}\to1$.
\end{enumerate}
\resultado (a) $\tfrac53$; (b) $1$; (c) $2$; (d) $\ln2\approx0{,}6931$; (e) $1$; (f) $1$.
\comprobacion numérica con $x=10^{-6}$: (a) $1{,}6666667$; (c) $2{,}000002$; (d) $0{,}6931474$.
\end{solucion}

\begin{solucion}{eje:2-asintotas}
\datos cuatro funciones.
\planteamiento verticales: puntos donde algún lateral es infinito; horizontales y oblicuas: \cref{prop:oblicua}.
\pasos
\begin{enumerate}[label=(\alph*)]
  \item División: $2x^{2}-3x=(x+1)(2x-5)+5$, luego $f(x)=2x-5+\frac{5}{x+1}$. Vertical $x=-1$ ($f\to\pm\infty$ por la derecha y la izquierda, respectivamente); oblicua $y=2x-5$ en $\pm\infty$.
  \item $g$ está definida en $\R$ y es continua: no hay verticales. Para $x>0$, $g(x)=\frac{1}{\sqrt{1+1/x^{2}}}\to1$; para $x<0$, $\sqrt{x^{2}}=-x$ y $g(x)=\frac{-1}{\sqrt{1+1/x^{2}}}\to-1$. Horizontales $y=1$ en $+\infty$ e $y=-1$ en $-\infty$.
  \item $h(x)=xe^{1/x}$, $x\ne0$. Si $x\to0^{+}$, con $t=1/x\to+\infty$, $h=e^{t}/t\to+\infty$ (\cref{prop:crecimientos}): vertical $x=0$ por la derecha. Si $x\to0^{-}$, $e^{1/x}\to0$ y $h\to0$. En $\pm\infty$: $m=\lim e^{1/x}=1$ y $b=\lim x(e^{1/x}-1)=\lim_{t\to0}\frac{e^{t}-1}{t}=1$. Oblicua $y=x+1$.
  \item $k(x)=\softplus(2x)$: en $-\infty$, $k\to\ln1=0$ (horizontal $y=0$); en $+\infty$, $k(x)-2x=\softplus(-2x)\to0$ (oblicua $y=2x$). No hay verticales.
\end{enumerate}
\resultado (a) $x=-1$, $y=2x-5$; (b) $y=1$ y $y=-1$; (c) $x=0$ (por la derecha), $y=x+1$; (d) $y=0$ en $-\infty$, $y=2x$ en $+\infty$.
\comprobacion (c) en $x=100$: $100e^{0{,}01}=101{,}005$, frente a $x+1=101$.
\end{solucion}

\begin{solucion}{eje:2-eN}
\datos cuatro sucesiones y sus límites.
\planteamiento trabajo previo: acotar $\abs{a_n-L}$ por $C/n$ o $C/\sqrt n$; después, redactar con los cuantificadores en orden.
\pasos
\begin{enumerate}[label=(\alph*)]
  \item $\abs{\frac{n^{2}}{n^{2}+1}-1}=\frac{1}{n^{2}+1}<\frac1n$. Sea $\varepsilon>0$ y $N\in\N$ con $N>1/\varepsilon$. Para $n\ge N$, $\abs{a_n-1}<\frac1n\le\frac1N<\varepsilon$.
  \item Sea $\varepsilon>0$ y $N\in\N$ con $N>1/\varepsilon^{2}$. Para $n\ge N$, $\sqrt n\ge\sqrt N>1/\varepsilon$, luego $\frac{1}{\sqrt n}<\varepsilon$.
  \item $\abs{\frac{(-1)^{n}}{n^{2}}}=\frac{1}{n^{2}}\le\frac1n$. Sea $\varepsilon>0$ y $N>1/\varepsilon$; para $n\ge N$, $\abs{a_n}\le\frac1n<\varepsilon$.
  \item $\abs{\frac{n+5}{2n-1}-\frac12}=\frac{\abs{2n+10-2n+1}}{2(2n-1)}=\frac{11}{2(2n-1)}\le\frac{11}{2n}$, porque $2n-1\ge n$ para $n\ge1$. Sea $\varepsilon>0$ y $N>\frac{11}{2\varepsilon}$; para $n\ge N$, $\abs{a_n-\frac12}\le\frac{11}{2n}<\varepsilon$. \hfill$\blacksquare$
\end{enumerate}
\resultado las cuatro convergencias, con $N>1/\varepsilon$, $N>1/\varepsilon^{2}$, $N>1/\varepsilon$ y $N>11/(2\varepsilon)$.
\comprobacion (d) con $\varepsilon=0{,}1$: $N=56$ y $\abs{a_{56}-\frac12}=\frac{11}{222}\approx0{,}0495<0{,}1$.
\end{solucion}

\begin{solucion}{eje:2-ed}
\datos cuatro límites.
\planteamiento factorizar $\abs{f(x)-L}$ como $\abs{x-a}$ por un factor y acotar ese factor (\cref{ej:cuadratico}).
\pasos
\begin{enumerate}[label=(\alph*)]
  \item $\abs{(5x+8)-18}=5\abs{x-2}$. Dado $\varepsilon>0$, sea $\delta=\varepsilon/5$. Si $0<\abs{x-2}<\delta$, $\abs{(5x+8)-18}<5\delta=\varepsilon$.
  \item $\abs{x^{2}+2x-8}=\abs{x-2}\abs{x+4}$. Si $\abs{x-2}<1$, entonces $1<x<3$ y $\abs{x+4}<7$. Dado $\varepsilon>0$, sea $\delta=\min(1,\varepsilon/7)$. Si $0<\abs{x-2}<\delta$, $\abs{x^{2}+2x-8}<7\abs{x-2}<7\cdot\frac{\varepsilon}{7}=\varepsilon$.
  \item Para $x\ne3$, $\frac{x^{2}-9}{x-3}=x+3$ y $\abs{(x+3)-6}=\abs{x-3}$. Dado $\varepsilon>0$, $\delta=\varepsilon$ sirve; la condición $0<\abs{x-3}$ es esencial, pues la función no está definida en~3.
  \item Dado $\varepsilon>0$, sea $\delta=\varepsilon^{1/4}$. Si $0<\abs{x}<\delta$, $\abs{x^{4}}=\abs{x}^{4}<\delta^{4}=\varepsilon$. \hfill$\blacksquare$
\end{enumerate}
\resultado $\delta=\varepsilon/5$; $\delta=\min(1,\varepsilon/7)$; $\delta=\varepsilon$; $\delta=\varepsilon^{1/4}$.
\comprobacion (b) con $\varepsilon=0{,}07$: $\delta=0{,}01$ y, en el peor caso $x=2{,}01$, $x^{2}+2x=8{,}0601$, a distancia $0{,}0601<0{,}07$.
\end{solucion}

\begin{solucion}{eje:2-ed-dificiles}
\datos tres límites.
\planteamiento como en el ejercicio anterior; en (b), acotar $1/\abs x$ alejando $x$ de~0; en (c), usar la \cref{def:limites-variantes}(ii).
\pasos
\begin{enumerate}[label=(\alph*)]
  \item Para $x\ge0$, $\abs{\sqrt x-2}=\frac{\abs{x-4}}{\sqrt x+2}\le\frac{\abs{x-4}}{2}$. Dado $\varepsilon>0$, sea $\delta=\min(4,2\varepsilon)$; la condición $\delta\le4$ garantiza $x>0$, de modo que $\sqrt x$ está definida. Si $0<\abs{x-4}<\delta$, $\abs{\sqrt x-2}\le\frac{\abs{x-4}}{2}<\varepsilon$.
  \item $\abs{\frac1x-1}=\frac{\abs{x-1}}{\abs x}$. Si $\abs{x-1}<\frac12$, entonces $x>\frac12$ y $\frac{1}{\abs x}<2$. Dado $\varepsilon>0$, sea $\delta=\min(\frac12,\frac\varepsilon2)$; si $0<\abs{x-1}<\delta$, $\abs{\frac1x-1}<2\abs{x-1}<\varepsilon$.
  \item Sea $K\in\R$. Si $K\le0$, $\frac{1}{x^{2}}>0\ge K$ para todo $x\ne0$ y cualquier $\delta$ sirve. Si $K>0$, sea $\delta=1/\sqrt K$; si $0<\abs x<\delta$, $x^{2}<1/K$ y $\frac{1}{x^{2}}>K$. \hfill$\blacksquare$
\end{enumerate}
\resultado $\delta=\min(4,2\varepsilon)$; $\delta=\min(\frac12,\frac\varepsilon2)$; $\delta=1/\sqrt K$.
\comprobacion (a) con $\varepsilon=0{,}1$: $\delta=0{,}2$; $\sqrt{4{,}2}=2{,}0494$ y $\sqrt{3{,}8}=1{,}9494$, ambos a distancia menor que $0{,}1$ de~2.
\end{solucion}

\begin{solucion}{eje:2-no-existe}
\datos $f(x)=\frac{\abs{x-1}}{x-1}$; la función de Dirichlet $D$.
\planteamiento (a) reducción al absurdo con la definición; (b) dos sucesiones con imágenes de límites distintos (\cref{teo:heine}).
\pasos
\begin{enumerate}[label=\arabic*.]
  \item (a) Supóngase $\lim_{x\to1}f(x)=L$. Con $\varepsilon=1$ existe $\delta>0$ tal que $\abs{f(x)-L}<1$ si $0<\abs{x-1}<\delta$.
  \item Para $x=1+\delta/2$, $f(x)=1$, luego $\abs{1-L}<1$ y $L>0$. Para $x=1-\delta/2$, $f(x)=-1$, luego $\abs{-1-L}<1$ y $L<0$. Contradicción: el límite no existe.
  \item (b) Sean $x_n=1/n\in\Q$ e $y_n=\sqrt2/n\notin\Q$ (si fuera racional, $\sqrt2=ny_n$ también lo sería). Ambas tienden a~0 sin anularse, y $D(x_n)=1\to1$, $D(y_n)=0\to0$. Por el \cref{teo:heine}, $\lim_{x\to0}D(x)$ no existe. \hfill$\blacksquare$
\end{enumerate}
\resultado ninguno de los dos límites existe.
\comprobacion en (a) los límites laterales existen y valen $1$ y $-1$; en (b) ni siquiera existen los laterales, porque el mismo argumento vale con sucesiones positivas o negativas.
\end{solucion}

\begin{solucion}{eje:2-producto}
\datos $f\to L$ y $g\to M$ cuando $x\to a$.
\planteamiento $f(x)g(x)-LM=f(x)\bigl(g(x)-M\bigr)+M\bigl(f(x)-L\bigr)$; acotar $\abs{f(x)}$ cerca de $a$.
\pasos
\begin{enumerate}[label=\arabic*.]
  \item Con $\varepsilon=1$ en la definición para $f$, existe $\delta_1>0$ tal que $\abs{f(x)-L}<1$, y por tanto $\abs{f(x)}<\abs L+1=:K$, si $0<\abs{x-a}<\delta_1$.
  \item Sea $\varepsilon>0$. Existe $\delta_2>0$ con $\abs{g(x)-M}<\frac{\varepsilon}{2K}$ si $0<\abs{x-a}<\delta_2$, y existe $\delta_3>0$ con $\abs{f(x)-L}<\frac{\varepsilon}{2(\abs M+1)}$ si $0<\abs{x-a}<\delta_3$.
  \item Sea $\delta=\min(\delta_1,\delta_2,\delta_3)$. Si $0<\abs{x-a}<\delta$,
  \begin{align*}
  \abs{f(x)g(x)-LM}&\le\abs{f(x)}\abs{g(x)-M}+\abs M\abs{f(x)-L}\\
  &<K\cdot\frac{\varepsilon}{2K}+\abs M\cdot\frac{\varepsilon}{2(\abs M+1)}<\frac\varepsilon2+\frac\varepsilon2=\varepsilon.
  \end{align*}
  Luego $\lim_{x\to a}f(x)g(x)=LM$. \hfill$\blacksquare$
\end{enumerate}
\resultado la ley del producto.
\comprobacion se divide por $\abs M+1$ y no por $\abs M$ para cubrir el caso $M=0$; con $f(x)=g(x)=x$ y $a=2$, la ley da $\lim x^{2}=4$, coherente con el \cref{ej:cuadratico}.
\end{solucion}

\begin{solucion}{eje:2-cociente}
\datos $b_n\to M\ne0$, $b_n\ne0$; la sucesión $(a_n)$ de reales no nulos.
\planteamiento (a) alejar $b_n$ de~0 y acotar $\abs{1/b_n-1/M}$; (b) contraejemplo en (i) y (a) en (ii).
\pasos
\begin{enumerate}[label=\arabic*.]
  \item Con $\varepsilon=\abs M/2$ existe $N_1$ tal que $\abs{b_n-M}<\abs M/2$ para $n\ge N_1$, y entonces $\abs{b_n}\ge\abs M-\abs{M-b_n}>\abs M/2$.
  \item Para $n\ge N_1$, $\abs{\frac1{b_n}-\frac1M}=\frac{\abs{M-b_n}}{\abs{b_n}\abs M}<\frac{2}{M^{2}}\abs{b_n-M}$.
  \item Sea $\varepsilon>0$. Existe $N_2$ con $\abs{b_n-M}<\frac{\varepsilon M^{2}}{2}$ para $n\ge N_2$. Para $n\ge\max(N_1,N_2)$, $\abs{\frac1{b_n}-\frac1M}<\varepsilon$. Esto prueba que $1/b_n\to1/M$, y con el \cref{teo:algebra-sucesiones}(iii), $a_n/b_n\to L/M$.
  \item (b)(i) No basta: $a_n=1/n\to0$, y $1/a_n=n$ no converge.
  \item (b)(ii) Basta: es exactamente el apartado (a). \hfill$\blacksquare$
\end{enumerate}
\resultado $1/b_n\to1/M$; la condición (i) no garantiza la convergencia de $(1/a_n)$ y la (ii) sí.
\comprobacion $b_n=2+1/n\to2$: $1/b_n=\frac{n}{2n+1}\to\frac12$, y la cota del paso~2 da $\abs{\frac{n}{2n+1}-\frac12}=\frac{1}{2(2n+1)}<\frac{2}{4}\cdot\frac1n$.
\end{solucion}

\begin{solucion}{eje:2-alternada}
\datos $a_n=(-1)^{n}$; una sucesión convergente cualquiera.
\planteamiento probar primero (b) y deducir (a) y el caso $r\le-1$ por contrarrecíproco.
\pasos
\begin{enumerate}[label=\arabic*.]
  \item (b) Sea $a_n\to L$ y $\varepsilon>0$. Existe $N$ con $\abs{a_n-L}<\varepsilon/2$ para $n\ge N$; entonces también $\abs{a_{n+1}-L}<\varepsilon/2$ y $\abs{a_{n+1}-a_n}\le\abs{a_{n+1}-L}+\abs{L-a_n}<\varepsilon$. Luego $a_{n+1}-a_n\to0$.
  \item (a) Para $a_n=(-1)^{n}$, $\abs{a_{n+1}-a_n}=2$ para todo $n$, que no tiende a~0. Por el contrarrecíproco de (b), $(a_n)$ diverge.
  \item Si $r\le-1$, $\abs{r^{n+1}-r^{n}}=\abs r^{n}\abs{r-1}\ge1\cdot2=2$, y $(r^{n})$ diverge por la misma razón.
  \item (c) El recíproco es falso: con $a_n=\sqrt n$, $a_{n+1}-a_n=\frac{1}{\sqrt{n+1}+\sqrt n}\to0$, pero $\sqrt n\to+\infty$. \hfill$\blacksquare$
\end{enumerate}
\resultado (a) y (b) demostrados; (c) el recíproco es falso.
\comprobacion $\sqrt{10^{6}+1}-\sqrt{10^{6}}\approx0{,}0005$, mientras que $\sqrt{10^{6}}=1000$.
\end{solucion}

\begin{solucion}{eje:2-recurrencia}
\datos $a_1=1$, $a_{n+1}=\sqrt{2+a_n}$.
\planteamiento inducción para acotación y monotonía; \cref{teo:monotona}; paso al límite en la recurrencia.
\pasos
\begin{enumerate}[label=\arabic*.]
  \item \emph{Acotación.} $0<a_1=1<2$. Si $0<a_n<2$, entonces $2<2+a_n<4$ y $\sqrt2<a_{n+1}<2$. Por inducción, $0<a_n<2$ para todo $n$.
  \item \emph{Monotonía.} Como $a_n>0$ y $a_{n+1}>0$, $a_{n+1}>a_n\iff2+a_n>a_n^{2}\iff(2-a_n)(1+a_n)>0$, que es cierto porque $0<a_n<2$. Luego $(a_n)$ es estrictamente creciente.
  \item Por el \cref{teo:monotona}, $a_n\to L$ con $1\le L\le2$.
  \item Elevando al cuadrado, $a_{n+1}^{2}=2+a_n$. Como $(a_{n+1})$ también tiende a $L$, el álgebra de límites da $L^{2}=2+L$, es decir, $(L-2)(L+1)=0$, y como $L\ge1$, $L=2$. \hfill$\blacksquare$
\end{enumerate}
\resultado $(a_n)$ es creciente, está en $(0,2)$ y converge a~2.
\comprobacion $a_2=\sqrt3\approx1{,}732$, $a_3\approx1{,}9319$, $a_4\approx1{,}9829$, $a_5\approx1{,}9957$.
\end{solucion}

\begin{solucion}{eje:2-raizn}
\datos $\sqrt[n]{n}=1+h_n$ con $h_n\ge0$ (porque $n\ge1$).
\planteamiento desarrollar $(1+h_n)^{n}$ por el binomio quedándose con un término; sándwich.
\pasos
\begin{enumerate}[label=\arabic*.]
  \item (a) Sea $n\ge2$. Todos los términos del binomio son $\ge0$, luego $n=(1+h_n)^{n}\ge\binom n2h_n^{2}=\frac{n(n-1)}{2}h_n^{2}$.
  \item Despejando, $0\le h_n\le\sqrt{\frac{2}{n-1}}$, y el miembro derecho tiende a~0. Por el sándwich, $h_n\to0$ y $\sqrt[n]n\to1$.
  \item (b) Si $a\ge1$: para $n\ge a$, $1\le\sqrt[n]a\le\sqrt[n]n$, y el sándwich da $\sqrt[n]a\to1$.
  \item Si $0<a<1$: $\sqrt[n]a=1/\sqrt[n]{1/a}$ con $1/a>1$; por el caso anterior y el \cref{teo:algebra-sucesiones}(iv), $\sqrt[n]a\to1$. \hfill$\blacksquare$
\end{enumerate}
\resultado $\sqrt[n]n\to1$ y $\sqrt[n]a\to1$ para $a>0$.
\comprobacion $\sqrt[100]{100}\approx1{,}0471$ y la cota da $h_{100}\le\sqrt{2/99}\approx0{,}142$.
\end{solucion}

\begin{solucion}{eje:2-laterales}
\datos la \cref{def:limite-funcion} y la \cref{def:limites-variantes}(i).
\planteamiento ($\Rightarrow$) el mismo $\delta$ sirve; ($\Leftarrow$) tomar el mínimo de los dos $\delta$.
\pasos
\begin{enumerate}[label=\arabic*.]
  \item ($\Rightarrow$) Sea $\varepsilon>0$ y $\delta$ el del límite. Si $a<x<a+\delta$, entonces $0<\abs{x-a}<\delta$ y $\abs{f(x)-L}<\varepsilon$; igual si $a-\delta<x<a$. Luego ambos laterales valen $L$.
  \item ($\Leftarrow$) Sea $\varepsilon>0$, y sean $\delta_{+}$ y $\delta_{-}$ los de los límites laterales. Con $\delta=\min(\delta_{+},\delta_{-})$, si $0<\abs{x-a}<\delta$, o bien $a<x<a+\delta_{+}$ o bien $a-\delta_{-}<x<a$; en ambos casos $\abs{f(x)-L}<\varepsilon$.
  \item Aplicación: $\frac{\ReLU(x)}{x}=1$ si $x>0$ y $0$ si $x<0$. Los laterales en~0 son $1$ y $0$, distintos, luego el límite no existe. \hfill$\blacksquare$
\end{enumerate}
\resultado la equivalencia; $\lim_{x\to0}\ReLU(x)/x$ no existe.
\comprobacion los laterales $1$ y $0$ son las derivadas laterales de la $\ReLU$ en~0 (\cref{ej:def-derivada}(d)).
\end{solucion}
% ---------------------------------------------------------------------
\section{Tema 3}

\begin{solucion}{eje:3-verdadero-falso}
\datos cuatro afirmaciones sobre funciones continuas.
\planteamiento buscar contraejemplos para las falsas y demostrar la verdadera.
\pasos
\begin{enumerate}[label=(\alph*)]
  \item \textbf{Falsa.} $\cos$ en $[-\pi,2\pi]$: $\cos(-\pi)=-1<0<1=\cos(2\pi)$, y $\cos$ se anula en $-\frac\pi2$, $\frac\pi2$ y $\frac{3\pi}{2}$: tres ceros. Bolzano garantiza al menos uno, no uno solo.
  \item \textbf{Falsa.} $(x-1)^{2}$ en $[-2,2]$: vale $9$ y $1$ en los extremos y se anula en $x=1$.
  \item \textbf{Falsa.} Sea $f(x)=1$ si $x\in\Q$ y $f(x)=-1$ si no. $\abs f=1$ es continua, pero $f$ no es continua en ningún punto (como la función de Dirichlet, \cref{eje:2-no-existe}).
  \item \textbf{Verdadera.} Por la desigualdad triangular inversa, $\bigl|\abs{f(x)}-\abs{f(a)}\bigr|\le\abs{f(x)-f(a)}$; el mismo $\delta$ que sirve para $f$ sirve para $\abs f$. (También: $\abs f$ es la composición de $f$ con $\abs{\cdot}$, \cref{prop:operaciones-continuas}.)
\end{enumerate}
\resultado (a) F; (b) F; (c) F; (d) V.
\comprobacion (c) y (d) muestran que la implicación «$f$ continua $\Rightarrow\abs f$ continua» no tiene recíproco.
\end{solucion}

\begin{solucion}{eje:3-hipotesis-weierstrass}
\datos el \cref{teo:weierstrass} exige $f$ continua en un intervalo cerrado y acotado.
\planteamiento en cada caso se elimina exactamente una hipótesis.
\pasos
\begin{enumerate}[label=(\alph*)]
  \item $f(x)=1/x$ en $(0,1]$: continua, no acotada superiormente; no tiene máximo. Falla que el intervalo sea cerrado.
  \item $f(x)=\arctan x$ en $[0,\infty)$: continua y acotada, con $\sup=\pi/2$ no alcanzado. Falla que el intervalo sea acotado.
  \item $f(x)=x$ en $[0,1)$ y $f(1)=0$: acotada, con $\sup=1$ no alcanzado; solo es discontinua en $x=1$. Falla la continuidad.
\end{enumerate}
\resultado los tres ejemplos anteriores.
\comprobacion en (b), $\arctan x<\pi/2$ para todo $x$, pero $\arctan x\to\pi/2$: el supremo es $\pi/2$ y no es máximo.
\end{solucion}

\begin{solucion}{eje:3-clasificar}
\datos cinco funciones.
\planteamiento por la \cref{prop:elementales-continuas}, cada función es continua en su dominio salvo, quizá, en los puntos de cambio de definición; en los puntos aislados donde no está definida se estudian los límites laterales y se clasifica según la \cref{def:discontinuidades}.
\pasos
\begin{enumerate}[label=(\alph*)]
  \item $f(x)=\frac{x-2}{(x-2)(x+2)}=\frac{1}{x+2}$ para $x\ne\pm2$. En $x=2$, el límite es $\frac14$: discontinuidad \emph{evitable}. En $x=-2$, los laterales son $\pm\infty$: \emph{esencial} (asíntota vertical).
  \item $1/\sin x$ no está definida en $x=k\pi$, $k\in\Z$, y allí $\sin x\to0$ con signos opuestos a cada lado: laterales infinitos, discontinuidad \emph{esencial}.
  \item $\floor x$ es constante en cada $[n,n+1)$, luego continua en los no enteros. En $n\in\Z$: $\lim_{x\to n^{-}}\floor x=n-1$ y $\lim_{x\to n^{+}}\floor x=n$: discontinuidad \emph{de salto}, de altura~1.
  \item Para $x\ne0$, $k$ es continua. Cuando $x\to0$, $-1/x^{2}\to-\infty$ y $e^{-1/x^{2}}\to0=k(0)$. $k$ es \textbf{continua en todo $\R$}.
  \item $\abs x/x$ vale $1$ para $x>0$ y $-1$ para $x<0$: en $x=0$, salto (no evitable, porque los laterales difieren).
\end{enumerate}
\resultado (a) evitable en $2$, esencial en $-2$; (b) esenciales en $k\pi$; (c) saltos en los enteros; (d) continua; (e) salto en~0.
\comprobacion (d) numérica: $k(0{,}1)=e^{-100}\approx3{,}7\cdot10^{-44}$.
\end{solucion}

\begin{solucion}{eje:3-constantes}
\datos cuatro funciones definidas a trozos.
\planteamiento en el punto de unión, igualar el valor de la función con los límites laterales.
\pasos
\begin{enumerate}[label=(\alph*)]
  \item Para $x\ne-2$, $\frac{(x+1)(x+2)}{x+2}=x+1\to-1$: $k=-1$.
  \item En $x=4$: el límite por la izquierda es $e^{4k}$ y por la derecha y el valor, $7$. $e^{4k}=7$, $k=\frac{\ln7}{4}\approx0{,}4865$.
  \item En $x=3$: $\sqrt{3k}=4$, $3k=16$, $k=\frac{16}{3}$.
  \item En $x=0$: $\lim_{x\to0^{+}}\frac{\sin2x}{x}=2\lim\frac{\sin2x}{2x}=2$; $\lim_{x\to0^{-}}(ax+b)=b$; $u(0)=a$. Continuidad: $a=b=2$.
\end{enumerate}
\resultado (a) $k=-1$; (b) $k=\frac{\ln7}{4}$; (c) $k=\frac{16}{3}$; (d) $a=b=2$.
\comprobacion (c): con $k=16/3$, $h(3)=\sqrt{16}=4=3+1$.
\end{solucion}

\begin{solucion}{eje:3-biseccion}
\datos $g(t)=\cos t-t^{3}$, continua en $\R$.
\planteamiento Bolzano en $[0,1]$ y bisección (\cref{teo:bolzano}).
\pasos
\begin{enumerate}[label=\arabic*.]
  \item (a) $g(0)=1>0$ y $g(1)=\cos1-1\approx-0{,}460<0$. Por Bolzano hay una solución en $(0,1)$.
  \item (b) $g(0{,}5)\approx0{,}753>0$: raíz en $[0{,}5;1]$. $g(0{,}75)\approx0{,}310>0$: en $[0{,}75;1]$. $g(0{,}875)\approx-0{,}029<0$: en $[0{,}75;0{,}875]$. $g(0{,}8125)\approx0{,}151>0$: en $[0{,}8125;0{,}875]$, de longitud $1/16$.
  \item (c) Tras $n$ pasos el intervalo mide $2^{-n}$. Para que cualquier punto del intervalo diste menos de $10^{-6}$ de la raíz basta $2^{-n}<10^{-6}$, es decir, $n\ge20$, porque $2^{20}=1\,048\,576$. Si se toma el punto medio, el error es menor que $2^{-(n+1)}$ y bastan $19$ pasos.
\end{enumerate}
\resultado (a) existe solución en $(0,1)$; (b) está en $[0{,}8125;\,0{,}875]$; (c) 20 pasos (19 con el punto medio).
\comprobacion la solución, calculada numéricamente, es $t\approx0{,}86547$, dentro del intervalo de (b).
\end{solucion}

\begin{solucion}{eje:3-raices}
\datos tres ecuaciones.
\planteamiento evaluar en enteros y aplicar Bolzano en cada cambio de signo.
\pasos
\begin{enumerate}[label=(\alph*)]
  \item $p(x)=x^{3}-3x+1$: $p(-2)=-1$, $p(-1)=3$, $p(0)=1$, $p(1)=-1$, $p(2)=3$. Cambios de signo en $(-2,-1)$, $(0,1)$ y $(1,2)$: tres raíces, y no hay más por ser de grado~3.
  \item $h(x)=e^{x}-3x$: $h(0)=1$, $h(1)=e-3\approx-0{,}28$, $h(2)=e^{2}-6\approx1{,}39$. Raíces en $(0,1)$ y $(1,2)$.
  \item $g(x)=\sigma(x)-x$: $g(0)=\frac12>0$ y $g(1)=\sigma(1)-1\approx-0{,}27<0$. Raíz en $(0,1)$.
\end{enumerate}
\resultado (a) en $(-2,-1)$, $(0,1)$, $(1,2)$; (b) en $(0,1)$ y $(1,2)$; (c) en $(0,1)$.
\comprobacion valores numéricos de las raíces: (a) $-1{,}879$; $0{,}347$; $1{,}532$; (b) $0{,}619$ y $1{,}512$; (c) $0{,}659$.
\end{solucion}

\begin{solucion}{eje:3-ed-continua}
\datos tres funciones.
\planteamiento \cref{def:continuidad}: dado $\varepsilon$, encontrar $\delta$ (que puede depender de $a$).
\pasos
\begin{enumerate}[label=(\alph*)]
  \item $\abs{(3x-7)-(3a-7)}=3\abs{x-a}$. Dado $\varepsilon>0$, $\delta=\varepsilon/3$.
  \item $\abs{x^{2}-a^{2}}=\abs{x-a}\abs{x+a}$, y $\abs{x+a}\le\abs{x-a}+2\abs a<1+2\abs a$ si $\abs{x-a}<1$. Dado $\varepsilon>0$, sea $\delta=\min\bigl(1,\frac{\varepsilon}{1+2\abs a}\bigr)$; si $\abs{x-a}<\delta$, $\abs{x^{2}-a^{2}}<(1+2\abs a)\delta\le\varepsilon$.
  \item Si $a>0$: para $x\ge0$, $\abs{\sqrt x-\sqrt a}=\frac{\abs{x-a}}{\sqrt x+\sqrt a}\le\frac{\abs{x-a}}{\sqrt a}$; dado $\varepsilon>0$, $\delta=\varepsilon\sqrt a$ sirve. Si $a=0$: para $x\ge0$, $\sqrt x<\varepsilon\iff x<\varepsilon^{2}$; $\delta=\varepsilon^{2}$ sirve. \hfill$\blacksquare$
\end{enumerate}
\resultado $\delta=\varepsilon/3$; $\delta=\min(1,\varepsilon/(1+2\abs a))$; $\delta=\varepsilon\sqrt a$ o $\delta=\varepsilon^{2}$.
\comprobacion en (b), $\delta$ disminuye al crecer $\abs a$: la parábola es más empinada lejos del origen, y el mismo $\varepsilon$ exige un $\delta$ menor. Ningún $\delta$ sirve a la vez para todos los $a$.
\end{solucion}

\begin{solucion}{eje:3-punto-fijo}
\datos $f\colon[0,1]\to[0,1]$ continua.
\planteamiento aplicar Bolzano a $g(x)=f(x)-x$.
\pasos
\begin{enumerate}[label=\arabic*.]
  \item $g$ es continua en $[0,1]$, $g(0)=f(0)\ge0$ y $g(1)=f(1)-1\le0$, porque $f$ toma valores en $[0,1]$.
  \item Si $g(0)=0$ o $g(1)=0$, el punto fijo es $0$ o $1$. Si no, $g(0)>0>g(1)$ y, por Bolzano (aplicado a $-g$), existe $c\in(0,1)$ con $g(c)=0$, es decir, $f(c)=c$.
  \item (b) $f(x)=x^{2}$ es continua de $(0,1)$ en $(0,1)$ y $x^{2}<x$ en $(0,1)$: no hay punto fijo. $f(x)=1$ si $x<\frac12$ y $f(x)=0$ si $x\ge\frac12$ va de $[0,1]$ en $[0,1]$ y no tiene punto fijo: $f(x)=x$ exigiría $x=1<\frac12$ o $x=0\ge\frac12$. \hfill$\blacksquare$
\end{enumerate}
\resultado existe $c$ con $f(c)=c$; los ejemplos de (b) muestran que hacen falta el intervalo cerrado y la continuidad.
\comprobacion con $f=\cos$, que lleva $[0,1]$ en $[\cos1,1]\subseteq[0,1]$, el punto fijo es $c\approx0{,}739$.
\end{solucion}

\begin{solucion}{eje:3-cuerda}
\datos posiciones $s_1(t)$ (subida) y $s_2(t)$ (bajada), medidas como distancia a la base a lo largo del camino, continuas en $t\in[8,14]$; longitud del camino $D>0$.
\planteamiento Bolzano aplicado a $h=s_1-s_2$.
\pasos
\begin{enumerate}[label=\arabic*.]
  \item $s_1(8)=0$, $s_1(14)=D$, $s_2(8)=D$, $s_2(14)=0$.
  \item $h=s_1-s_2$ es continua, $h(8)=-D<0$ y $h(14)=D>0$.
  \item Por Bolzano, existe $t_0\in(8,14)$ con $h(t_0)=0$: a la hora $t_0$, el ciclista está el primer día y el segundo en el mismo punto del camino. \hfill$\blacksquare$
\end{enumerate}
\resultado existe tal hora y tal punto.
\comprobacion el argumento no dice cuál es la hora ni que sea única; si un día se detiene a descansar, puede haber varias.
\end{solucion}

\begin{solucion}{eje:3-sucesiones}
\datos $f(x_n)\to f(y)$ para toda $(x_n)\subseteq[a,b]$ con $x_n\to y$.
\planteamiento contrarrecíproco: si $f$ no es continua en $y$, construir una sucesión mala.
\pasos
\begin{enumerate}[label=\arabic*.]
  \item Supóngase que $f$ no es continua en $y$. Negando la \cref{def:continuidad}: existe $\varepsilon>0$ tal que para todo $\delta>0$ hay $x\in[a,b]$ con $\abs{x-y}<\delta$ y $\abs{f(x)-f(y)}\ge\varepsilon$.
  \item Para cada $n\in\N$, con $\delta=1/n$, se elige $x_n\in[a,b]$ con $\abs{x_n-y}<1/n$ y $\abs{f(x_n)-f(y)}\ge\varepsilon$.
  \item Entonces $x_n\to y$, pero $f(x_n)\not\to f(y)$, en contra de la hipótesis. Luego $f$ es continua en $y$. \hfill$\blacksquare$
\end{enumerate}
\resultado $f$ es continua en $y$.
\comprobacion el recíproco (continua implica que conserva los límites de sucesiones) es la parte ($\Rightarrow$) del \cref{teo:heine}; juntos dan la caracterización por sucesiones de la continuidad.
\end{solucion}

\begin{solucion}{eje:3-inyectiva-monotona}
\datos $f$ continua e inyectiva en $[a,b]$, $f(a)<f(b)$.
\planteamiento (a) reducción al absurdo con el TVI; (b) aplicar (a) a subintervalos.
\pasos
\begin{enumerate}[label=\arabic*.]
  \item (a) Sea $x\in(a,b)$. Por la inyectividad, $f(x)\ne f(a)$ y $f(x)\ne f(b)$.
  \item Si fuera $f(x)>f(b)$, entonces $f(a)<f(b)<f(x)$ y el \cref{cor:tvi} en $[a,x]$ daría $c\in(a,x)$ con $f(c)=f(b)$; como $c<x<b$, $c\ne b$, en contra de la inyectividad.
  \item Si fuera $f(x)<f(a)$, entonces $f(x)<f(a)<f(b)$ y el TVI en $[x,b]$ daría $c\in(x,b)$ con $f(c)=f(a)$, $c\ne a$, de nuevo contradicción. Luego $f(a)<f(x)<f(b)$.
  \item (b) Sean $a\le x<y\le b$. Por (a) (o por la hipótesis si $x=a$), $f(x)<f(b)$ cuando $x<b$. La restricción de $f$ a $[x,b]$ es continua, inyectiva y cumple $f(x)<f(b)$; si $y<b$, (a) aplicado a ella da $f(x)<f(y)$, y si $y=b$ ya se sabe. Luego $f$ es estrictamente creciente. \hfill$\blacksquare$
\end{enumerate}
\resultado (a) $f(a)<f(x)<f(b)$ en $(a,b)$; (b) $f$ es estrictamente creciente.
\comprobacion la continuidad es imprescindible: la función inyectiva y no monótona del \cref{eje:1-estricta} tiene un salto.
\end{solucion}

\begin{solucion}{eje:3-cauchy}
\datos $g(x+y)=g(x)+g(y)$ para todos $x,y$.
\planteamiento (a) sustituciones particulares e inducción; (b) densidad de $\Q$ y continuidad; (c) trasladar la continuidad con la ecuación.
\pasos
\begin{enumerate}[label=\arabic*.]
  \item Con $x=y=0$: $g(0)=2g(0)$, luego $g(0)=0$. Con $y=-x$: $0=g(0)=g(x)+g(-x)$, luego $g(-x)=-g(x)$.
  \item $g(1\cdot x)=g(x)$, y si $g(nx)=ng(x)$, entonces $g((n+1)x)=g(nx)+g(x)=(n+1)g(x)$. Por inducción, vale para todo $n\in\N$.
  \item Sea $q=m/n$ con $m\in\Z$, $n\in\N$, y $c=g(1)$. Por el paso~2 (y la imparidad si $m<0$), $ng(m/n)=g(m)=mg(1)=mc$, luego $g(q)=cq$.
  \item (b) Sea $x\in\R$ y $(q_k)$ una sucesión de racionales con $q_k\to x$ (por ejemplo, los truncamientos decimales de $x$). Por la continuidad, $g(x)=\lim g(q_k)=\lim cq_k=cx$.
  \item (c) Sea $g$ continua en $x_0$ y sea $x\in\R$. Para todo $h$, $g(x+h)-g(x)=g(h)=g(x_0+h)-g(x_0)$, y el último miembro tiende a~0 cuando $h\to0$. Luego $g$ es continua en $x$. \hfill$\blacksquare$
\end{enumerate}
\resultado $g(q)=g(1)q$ en $\Q$; si $g$ es continua (basta en un punto), $g(x)=cx$ en $\R$.
\comprobacion $g(x)=cx$ cumple la ecuación: $c(x+y)=cx+cy$. Sin ninguna hipótesis de continuidad existen soluciones no lineales, pero su construcción requiere el axioma de elección.
\end{solucion}

\begin{solucion}{eje:3-exponencial-funcional}
\datos $h$ continua con $h(x+y)=h(x)h(y)$.
\planteamiento (a) escribir $h(x)$ como un cuadrado; (b) tomar logaritmos y aplicar el \cref{eje:3-cauchy}.
\pasos
\begin{enumerate}[label=\arabic*.]
  \item (a) $h(x)=h(\frac x2+\frac x2)=h(\frac x2)^{2}\ge0$ para todo $x$.
  \item Si $h(x_0)=0$ para algún $x_0$, entonces $h(x)=h(x_0)h(x-x_0)=0$ para todo $x$. Luego, o $h\equiv0$, o $h(x)\ne0$ para todo $x$, y entonces $h(x)>0$ por el paso~1.
  \item (b) Si $h>0$, $g=\ln h$ es continua (composición) y $g(x+y)=\ln\bigl(h(x)h(y)\bigr)=g(x)+g(y)$.
  \item Por el \cref{eje:3-cauchy}(b), $g(x)=cx$ con $c=g(1)$, luego $h(x)=e^{cx}=a^{x}$ con $a=e^{c}=h(1)>0$. \hfill$\blacksquare$
\end{enumerate}
\resultado $h\equiv0$ o $h(x)=a^{x}$ con $a=h(1)>0$.
\comprobacion $a^{x+y}=a^{x}a^{y}$; y $h(0)=h(0)^{2}$ da $h(0)\in\{0,1\}$, coherente con $a^{0}=1$.
\end{solucion}

\begin{solucion}{eje:3-periodica}
\datos $f$ continua en $\R$, $f(x+T)=f(x)$.
\planteamiento (a) reducir a $[0,T]$ y aplicar Weierstrass; (b) Bolzano a $g$.
\pasos
\begin{enumerate}[label=\arabic*.]
  \item (a) Por inducción, $f(x+kT)=f(x)$ para $k\in\N$, y también para $k$ negativo, porque $f(x)=f((x-T)+T)=f(x-T)$.
  \item Todo $x\in\R$ se escribe $x=kT+r$ con $k=\floor{x/T}\in\Z$ y $r\in[0,T)$, y $f(x)=f(r)$. Luego $f(\R)=f([0,T])$.
  \item Por Weierstrass, $f$ alcanza en $[0,T]$ un máximo y un mínimo, que lo son también en $\R$ por el paso~2.
  \item (b) Sea $g(x)=f(x+T/2)-f(x)$, continua. $g(0)=f(T/2)-f(0)$ y $g(T/2)=f(T)-f(T/2)=f(0)-f(T/2)=-g(0)$.
  \item Si $g(0)=0$, sirve $x=0$. Si no, $g(0)$ y $g(T/2)$ tienen signos opuestos y Bolzano da $x\in(0,T/2)$ con $g(x)=0$. \hfill$\blacksquare$
\end{enumerate}
\resultado $f$ alcanza máximo y mínimo; existe $x$ con $f(x+T/2)=f(x)$.
\comprobacion con $f=\sin$, $T=2\pi$: $\sin(x+\pi)=\sin x$ equivale a $-\sin x=\sin x$, que se cumple en $x=0$. Interpretación clásica: en todo momento hay dos puntos antípodas del ecuador con la misma temperatura.
\end{solucion}
% ---------------------------------------------------------------------
\section{Tema 4}

\begin{solucion}{eje:4-grafica-derivada}
\datos cuatro funciones; la derivada es la pendiente de la tangente.
\planteamiento leer la pendiente en cada tramo de la gráfica.
\pasos
\begin{enumerate}[label=(\alph*)]
  \item Constante: pendiente nula; $f'=0$ (el eje de abscisas).
  \item $\ReLU$: pendiente $0$ para $x<0$ y $1$ para $x>0$; $\ReLU'$ es la función escalón, no definida en~0.
  \item $\abs x$: pendiente $-1$ para $x<0$ y $1$ para $x>0$; la derivada es la función signo, no definida en~0.
  \item $\sigma$: siempre creciente, con la máxima pendiente en el centro de simetría; $\sigma'$ es una «campana» positiva, par, con máximo $\sigma'(0)=\frac14$, que tiende a~0 en $\pm\infty$.
  \item Paridad: si la gráfica es simétrica respecto del eje de ordenadas, la tangente en $-x$ es la reflejada de la tangente en $x$, y una reflexión vertical cambia el signo de la pendiente: $f'(-x)=-f'(x)$.
\end{enumerate}
\resultado (a) $0$; (b) escalón; (c) signo; (d) campana con máximo $\frac14$ en $x=0$; la derivada de una función par es impar.
\comprobacion (c) es la derivada de una función par ($\abs x$) y es impar; (d) es la derivada de $\sigma-\frac12$, impar, y es par (\cref{eje:4-par-impar}).
\end{solucion}

\begin{solucion}{eje:4-definicion}
\datos $f(x)=x^{3}$, $\sqrt x$ y $\frac{1}{2x+1}$.
\planteamiento cociente incremental, simplificación del factor $h$ y límite.
\pasos
\begin{enumerate}[label=(\alph*)]
  \item $\frac{(a+h)^{3}-a^{3}}{h}=\frac{3a^{2}h+3ah^{2}+h^{3}}{h}=3a^{2}+3ah+h^{2}\to3a^{2}$.
  \item $\frac{\sqrt{a+h}-\sqrt a}{h}=\frac{(a+h)-a}{h(\sqrt{a+h}+\sqrt a)}=\frac{1}{\sqrt{a+h}+\sqrt a}\to\frac{1}{2\sqrt a}$, por la continuidad de la raíz.
  \item $f(1)=\frac13$ y $\frac{1}{h}\Bigl(\frac{1}{3+2h}-\frac13\Bigr)=\frac{3-(3+2h)}{3h(3+2h)}=\frac{-2}{3(3+2h)}\to-\frac29$. Tangente: $y=\frac13-\frac29(x-1)$, es decir, $y=\frac{5-2x}{9}$.
\end{enumerate}
\resultado (a) $3a^{2}$; (b) $\frac{1}{2\sqrt a}$; (c) $f'(1)=-\frac29$, tangente $y=\frac{5-2x}{9}$.
\comprobacion (c) con la regla de la cadena: $f'(x)=-\frac{2}{(2x+1)^{2}}$, y $f'(1)=-\frac29$.
\end{solucion}

\begin{solucion}{eje:4-mecanico-1}
\datos seis funciones.
\planteamiento reglas de la suma, del producto, del cociente y de la cadena (\cref{teo:reglas,teo:cadena}).
\pasos
\begin{enumerate}[label=(\alph*)]
  \item $10x^{9}+15x^{4}$ ($\pi^{3}$ es constante).
  \item $(3x^{2}+1)(x^{5}+x^{2})+(x^{3}+x)(5x^{4}+2x)=8x^{7}+6x^{5}+5x^{4}+3x^{2}$.
  \item $\frac{(x^{2}-1)-(x+2)\cdot2x}{(x^{2}-1)^{2}}=-\frac{x^{2}+4x+1}{(x^{2}-1)^{2}}$.
  \item $\frac{x^{4}+1}{x}=x^{3}+x^{-1}$, luego $3x^{2}-\frac{1}{x^{2}}$.
  \item $100(x^{2}+2)^{99}\cdot2x=200x(x^{2}+2)^{99}$.
  \item $10x^{9}(x^{2}+1)^{10}+x^{10}\cdot10(x^{2}+1)^{9}\cdot2x=10x^{9}(x^{2}+1)^{9}(3x^{2}+1)$.
\end{enumerate}
\resultado los indicados.
\comprobacion (b) multiplicando primero: $(x^{3}+x)(x^{5}+x^{2})=x^{8}+x^{6}+x^{5}+x^{3}$, cuya derivada es $8x^{7}+6x^{5}+5x^{4}+3x^{2}$.
\end{solucion}

\begin{solucion}{eje:4-mecanico-2}
\datos ocho funciones.
\planteamiento regla de la cadena, \cref{prop:tabla,prop:derivadas-activaciones}.
\pasos
\begin{enumerate}[label=(\alph*)]
  \item $e^{-x^{2}/2}\cdot(-x)=-xe^{-x^{2}/2}$.
  \item $\frac{-\sin x}{\cos x}=-\tan x$.
  \item $2\sin3x\cdot\cos3x\cdot3=3\sin6x$.
  \item $\frac{1}{1+x}\cdot\frac{1}{2\sqrt x}=\frac{1}{2\sqrt x(1+x)}$.
  \item $\frac{1}{x+\sqrt{x^{2}+1}}\Bigl(1+\frac{x}{\sqrt{x^{2}+1}}\Bigr)=\frac{1}{x+\sqrt{x^{2}+1}}\cdot\frac{\sqrt{x^{2}+1}+x}{\sqrt{x^{2}+1}}=\frac{1}{\sqrt{x^{2}+1}}$.
  \item Con $z=wx+b$: $\frac{\partial}{\partial x}\sigma(z)=w\,\sigma(z)(1-\sigma(z))$ y $\frac{\partial}{\partial w}\sigma(z)=x\,\sigma(z)(1-\sigma(z))$.
  \item $2x\bigl(1-\tanh^{2}(x^{2})\bigr)$.
  \item $\frac{\sigma'(x)}{\sigma(x)}=1-\sigma(x)=\sigma(-x)$.
\end{enumerate}
\resultado los indicados.
\comprobacion (e) es la derivada de $\arsinh$ (\cref{eje:1-arsinh}), y coincide con la que da el \cref{teo:inversa}: $1/\cosh(\arsinh x)=1/\sqrt{1+\sinh^{2}(\arsinh x)}=1/\sqrt{1+x^{2}}$. (h) es coherente con $\ln\sigma(x)=-\softplus(-x)$ y $\softplus'=\sigma$.
\end{solucion}

\begin{solucion}{eje:4-logaritmica}
\datos tres funciones positivas.
\planteamiento derivar $\ln f$ y multiplicar por $f$ (\cref{ej:logaritmica}).
\pasos
\begin{enumerate}[label=(\alph*)]
  \item $\ln f=\sin x\ln x$; $\frac{f'}{f}=\cos x\ln x+\frac{\sin x}{x}$; $f'=x^{\sin x}\bigl(\cos x\ln x+\frac{\sin x}{x}\bigr)$.
  \item $\ln f=x\bigl(\ln(x+1)-\ln x\bigr)$; $\frac{f'}{f}=\ln\frac{x+1}{x}+x\Bigl(\frac{1}{x+1}-\frac1x\Bigr)=\ln\Bigl(1+\frac1x\Bigr)-\frac{1}{x+1}$.
  \item $\ln f=3\ln(x+1)+\frac12\ln(x^{2}+4)-5\ln(2x+3)$; $\frac{f'}{f}=\frac{3}{x+1}+\frac{x}{x^{2}+4}-\frac{10}{2x+3}$, y $f'=f\cdot\frac{f'}{f}$.
\end{enumerate}
\resultado (a) $x^{\sin x}\bigl(\cos x\ln x+\frac{\sin x}{x}\bigr)$; (b) $\bigl(1+\frac1x\bigr)^{x}\bigl(\ln(1+\frac1x)-\frac{1}{x+1}\bigr)$; (c) $\frac{(x+1)^{3}\sqrt{x^{2}+4}}{(2x+3)^{5}}\bigl(\frac{3}{x+1}+\frac{x}{x^{2}+4}-\frac{10}{2x+3}\bigr)$.
\comprobacion (b): por la \cref{prop:notables-exp}, $(1+1/x)^{x}\to e$; la derivada debe tender a~0, y en efecto $\ln(1+\frac1x)-\frac{1}{x+1}=\frac1x-\frac{1}{x+1}+O(x^{-2})\to0$.
\end{solucion}

\begin{solucion}{eje:4-implicita}
\datos (a) $x^{2/3}+y^{2/3}=4$, punto $(-\sqrt{27},1)$; (b) $\sin x+\sin y=\frac12$.
\planteamiento derivación implícita y despeje de $y'$.
\pasos
\begin{enumerate}[label=\arabic*.]
  \item (a) El punto está en la curva: $(-\sqrt{27})^{2/3}=27^{1/3}=3$ y $1^{2/3}=1$.
  \item Derivando: $\frac23x^{-1/3}+\frac23y^{-1/3}y'=0$, luego $y'=-\bigl(\frac yx\bigr)^{1/3}$.
  \item En el punto: $\frac yx=-\frac{1}{\sqrt{27}}=-3^{-3/2}$, su raíz cúbica es $-3^{-1/2}$, y $y'=3^{-1/2}=\frac{1}{\sqrt3}$.
  \item Tangente: $y=1+\frac{1}{\sqrt3}(x+\sqrt{27})=\frac{x}{\sqrt3}+4$, porque $\sqrt{27}/\sqrt3=3$.
  \item (b) Derivando: $\cos x+\cos y\cdot y'=0$, $y'=-\frac{\cos x}{\cos y}$ donde $\cos y\ne0$. Tangente horizontal: $\cos x=0$, es decir, $x=\frac\pi2+k\pi$, con $\sin x=\pm1$.
  \item Si $\sin x=1$: $\sin y=-\frac12$, $y=-\frac\pi6+2m\pi$ o $y=\frac{7\pi}{6}+2m\pi$, y ahí $\cos y=\pm\frac{\sqrt3}{2}\ne0$. Si $\sin x=-1$: $\sin y=\frac32$, imposible.
\end{enumerate}
\resultado (a) $y=\frac{x}{\sqrt3}+4$; (b) los puntos $\bigl(\frac\pi2+2k\pi,\,-\frac\pi6+2m\pi\bigr)$ y $\bigl(\frac\pi2+2k\pi,\,\frac{7\pi}{6}+2m\pi\bigr)$, con $k,m\in\Z$.
\comprobacion (a): la tangente pasa por el punto, $\frac{-\sqrt{27}}{\sqrt3}+4=-3+4=1$. (b): $\sin\frac\pi2+\sin(-\frac\pi6)=1-\frac12=\frac12$.
\end{solucion}

\begin{solucion}{eje:4-tangentes-origen}
\datos $f(x)=1+(x-1)^{2}$, $f'(x)=2(x-1)$.
\planteamiento tangente en un punto genérico $t$ e imposición de que pase por $(0,0)$.
\pasos
\begin{enumerate}[label=\arabic*.]
  \item Tangente en $x=t$: $y=1+(t-1)^{2}+2(t-1)(x-t)$.
  \item Pasa por el origen si $0=1+(t-1)^{2}-2t(t-1)=2-t^{2}$, es decir, $t=\pm\sqrt2$.
  \item En el origen la ordenada es $0$, luego la tangente es $y=f'(t)x=2(t-1)x$: para $t=\sqrt2$, $y=2(\sqrt2-1)x$; para $t=-\sqrt2$, $y=-2(\sqrt2+1)x$.
\end{enumerate}
\resultado $y=2(\sqrt2-1)x$ e $y=-2(\sqrt2+1)x$.
\comprobacion en $t=\sqrt2$: $f(\sqrt2)=1+(\sqrt2-1)^{2}=4-2\sqrt2$ y $2(\sqrt2-1)\sqrt2=4-2\sqrt2$.
\end{solucion}

\begin{solucion}{eje:4-pegado}
\datos $f$ definida a trozos con $x^{2}+4x+1$ y $ax+b$.
\planteamiento derivable implica continua (\cref{teo:derivable-continua}): igualar valores y derivadas laterales en el punto de unión.
\pasos
\begin{enumerate}[label=\arabic*.]
  \item (a) En $x=0$: continuidad, $b=0^{2}+0+1=1$; derivadas laterales, $a=2\cdot0+4=4$.
  \item Con $a=4$, $b=1$, ambos trozos tienen en~0 el mismo valor y la misma pendiente, luego $f'_{-}(0)=f'_{+}(0)=4$ y $f$ es derivable; en los demás puntos lo es por ser polinómica en cada trozo.
  \item (b) En $x=1$: continuidad, $a+b=1+4+1=6$; derivadas, $a=2+4=6$. Luego $a=6$, $b=0$.
\end{enumerate}
\resultado (a) $a=4$, $b=1$; (b) $a=6$, $b=0$.
\comprobacion en ambos casos $ax+b$ es la recta tangente a la parábola en el punto de unión: $y=1+4x$ en $x=0$ e $y=6+6(x-1)=6x$ en $x=1$.
\end{solucion}

\begin{solucion}{eje:4-sucesivas}
\datos $f(x)=xe^{x}$; $\sigma'=\sigma(1-\sigma)$; $\tanh'=1-\tanh^{2}$.
\planteamiento derivar repetidamente y detectar el patrón; inducción.
\pasos
\begin{enumerate}[label=\arabic*.]
  \item (a) $f'=(x+1)e^{x}$, $f''=(x+2)e^{x}$, y en general se conjetura $f^{(n)}=(x+n)e^{x}$. Si vale para $n$, $f^{(n+1)}=e^{x}+(x+n)e^{x}=(x+n+1)e^{x}$; vale para $n=1$, luego para todo $n$. En particular, $f^{(5)}(x)=(x+5)e^{x}$.
  \item (b) Con $s=\sigma(x)$: $\sigma''=\frac{\dd}{\dd x}(s-s^{2})=(1-2s)s'=s(1-s)(1-2s)=s-3s^{2}+2s^{3}$.
  \item $\sigma'''=(1-6s+6s^{2})s'$. En $x=0$, $s=\frac12$ y $s'=\frac14$: $\sigma'''(0)=(1-3+\frac32)\cdot\frac14=-\frac18$.
  \item (c) Con $t=\tanh x$: $\tanh''=\frac{\dd}{\dd x}(1-t^{2})=-2tt'=-2t(1-t^{2})$.
\end{enumerate}
\resultado (a) $f^{(n)}=(x+n)e^{x}$; (b) $\sigma''=\sigma(1-\sigma)(1-2\sigma)$, $\sigma'''(0)=-\frac18$; (c) $\tanh''=-2t(1-t^{2})$.
\comprobacion por el \cref{eje:1-sigmoide-simetria}, $\sigma(x)=\frac12+\frac12\tanh\frac x2$, luego $\sigma'''(0)=\frac12\cdot\frac18\tanh'''(0)$; y $\tanh'''=-2(1-t^{2})(1-3t^{2})$ vale $-2$ en~0, lo que da $-\frac18$.
\end{solucion}

\begin{solucion}{eje:4-producto}
\datos $f$, $g$ derivables en $a$.
\planteamiento sumar y restar $f(a+h)g(a)$; para Leibniz, inducción sobre $n$ con la fórmula de Pascal.
\pasos
\begin{enumerate}[label=\arabic*.]
  \item \emph{Enunciado.} Si $f$ y $g$ son derivables en $a$, $fg$ es derivable en $a$ y $(fg)'(a)=f'(a)g(a)+f(a)g'(a)$.
  \item Para $h\ne0$ pequeño, $\frac{f(a+h)g(a+h)-f(a)g(a)}{h}=f(a+h)\frac{g(a+h)-g(a)}{h}+g(a)\frac{f(a+h)-f(a)}{h}$ (se ha sumado y restado $f(a+h)g(a)$).
  \item Cuando $h\to0$: $\frac{g(a+h)-g(a)}{h}\to g'(a)$ y $\frac{f(a+h)-f(a)}{h}\to f'(a)$ por definición; $f(a+h)\to f(a)$ porque $f$, derivable en $a$, es continua en $a$ (\cref{teo:derivable-continua}). Este es el único lugar donde se usa la continuidad.
  \item Por el álgebra de límites, el cociente tiende a $f(a)g'(a)+g(a)f'(a)$.
  \item \emph{Leibniz.} Para $n=1$ es la regla del producto. Supóngase cierta para $n$ (con $f,g$ derivables $n+1$ veces). Derivando cada término con la regla del producto,
  \[
  (fg)^{(n+1)}=\sum_{k=0}^{n}\binom nk\bigl(f^{(k+1)}g^{(n-k)}+f^{(k)}g^{(n+1-k)}\bigr)=\sum_{j=0}^{n+1}\Bigl[\binom n{j-1}+\binom nj\Bigr]f^{(j)}g^{(n+1-j)},
  \]
  con el convenio $\binom n{-1}=\binom n{n+1}=0$, tras escribir $j=k+1$ en la primera suma y $j=k$ en la segunda.
  \item Por la fórmula de Pascal, $\binom n{j-1}+\binom nj=\binom{n+1}{j}$, que es la fórmula para $n+1$. Para $n=2$ resulta $f''g+2f'g'+fg''$. \hfill$\blacksquare$
\end{enumerate}
\resultado la regla del producto y la fórmula de Leibniz.
\comprobacion con $f=x$, $g=e^{x}$: $(xe^{x})^{(n)}=\binom n0xe^{x}+\binom n1\cdot1\cdot e^{x}=(x+n)e^{x}$, como en el \cref{eje:4-sucesivas}.
\end{solucion}

\begin{solucion}{eje:4-par-impar}
\datos $f$ derivable y par (o impar).
\planteamiento derivar la identidad $f(-x)=\pm f(x)$ con la regla de la cadena.
\pasos
\begin{enumerate}[label=\arabic*.]
  \item Si $f$ es par, $f(-x)=f(x)$ para todo $x$. Ambos miembros son funciones derivables de $x$; derivando el izquierdo con la regla de la cadena, $-f'(-x)=f'(x)$, luego $f'(-x)=-f'(x)$: $f'$ es impar.
  \item Si $f$ es impar, $f(-x)=-f(x)$; derivando, $-f'(-x)=-f'(x)$, luego $f'(-x)=f'(x)$: $f'$ es par.
  \item Comprobaciones: $\cosh$ es par y $\cosh'=\sinh$ es impar; $\sinh$ es impar y $\sinh'=\cosh$ es par; $\sigma-\frac12$ es impar (\cref{prop:sigmoide}(ii)) y su derivada $\sigma'(x)=\sigma(x)\sigma(-x)$ es par. \hfill$\blacksquare$
\end{enumerate}
\resultado par $\Rightarrow$ derivada impar; impar $\Rightarrow$ derivada par.
\comprobacion consecuencia: si $f$ es par y derivable en~0, $f'(0)=-f'(0)$, luego $f'(0)=0$ (el vértice de una función par derivable es un punto crítico).
\end{solucion}

\begin{solucion}{eje:4-funcional}
\datos $f(x+y)=f(x)+f(y)+x^{2}y+xy^{2}$; $f(x)/x\to1$ cuando $x\to0$.
\planteamiento sustituir valores particulares y escribir el cociente incremental con la ecuación.
\pasos
\begin{enumerate}[label=\arabic*.]
  \item Con $x=y=0$: $f(0)=2f(0)$, luego $f(0)=0$.
  \item $f'(0)=\lim_{h\to0}\frac{f(h)-f(0)}{h}=\lim\frac{f(h)}{h}=1$.
  \item Sea $x\in\R$. Con $y=h$: $f(x+h)-f(x)=f(h)+x^{2}h+xh^{2}$, luego $\frac{f(x+h)-f(x)}{h}=\frac{f(h)}{h}+x^{2}+xh\to1+x^{2}$.
  \item Por tanto $f$ es derivable en todo punto, con $f'(x)=1+x^{2}$. \hfill$\blacksquare$
\end{enumerate}
\resultado $f(0)=0$, $f'(0)=1$, $f'(x)=1+x^{2}$.
\comprobacion la función $f(x)=x+\frac{x^{3}}{3}$ cumple la ecuación: $\frac{(x+y)^{3}-x^{3}-y^{3}}{3}=x^{2}y+xy^{2}$. Por el \cref{cor:tvm}(ii), es la única solución, pues cualquier otra difiere de ella en una constante, que es nula porque $f(0)=0$.
\end{solucion}

\begin{solucion}{eje:4-raiz-doble}
\datos $p$ polinomio, $a\in\R$.
\planteamiento teorema del factor: $p(a)=0$ si y solo si $(x-a)$ divide a $p$.
\pasos
\begin{enumerate}[label=\arabic*.]
  \item ($\Rightarrow$) Si $p=(x-a)^{2}q$, entonces $p(a)=0$ y $p'=2(x-a)q+(x-a)^{2}q'$, luego $p'(a)=0$.
  \item ($\Leftarrow$) Si $p(a)=0$, $p=(x-a)q$ con $q$ polinomio. Derivando, $p'=q+(x-a)q'$, y $p'(a)=q(a)$. Si además $p'(a)=0$, entonces $q(a)=0$ y $q=(x-a)r$, de modo que $p=(x-a)^{2}r$.
  \item Aplicación: $p(x)=x^{3}-3x+c$, $p'(x)=3x^{2}-3=0\iff x=\pm1$. $p(1)=c-2=0$ da $c=2$; $p(-1)=c+2=0$ da $c=-2$. \hfill$\blacksquare$
\end{enumerate}
\resultado $a$ es raíz doble si y solo si $p(a)=p'(a)=0$; $c=2$ (raíz doble $1$) o $c=-2$ (raíz doble $-1$).
\comprobacion $x^{3}-3x+2=(x-1)^{2}(x+2)$ y $x^{3}-3x-2=(x+1)^{2}(x-2)$.
\end{solucion}

\begin{solucion}{eje:4-x2sen}
\datos $f(x)=x^{2}\sin(1/x)$, $g(x)=x\sin(1/x)$, ambas nulas en~0.
\planteamiento fuera de~0, reglas de derivación; en~0, la definición; para la discontinuidad de $f'$, una sucesión.
\pasos
\begin{enumerate}[label=\arabic*.]
  \item (a) Para $x\ne0$: $f'(x)=2x\sin\frac1x+x^{2}\cos\frac1x\cdot\bigl(-\frac{1}{x^{2}}\bigr)=2x\sin\frac1x-\cos\frac1x$.
  \item En~0: $\frac{f(h)-f(0)}{h}=h\sin\frac1h$, y $\abs{h\sin\frac1h}\le\abs h\to0$. Luego $f'(0)=0$ y $f$ es derivable en todo $\R$.
  \item (b) Sea $x_n=\frac{1}{2\pi n}\to0$. Entonces $f'(x_n)=0-\cos(2\pi n)=-1$, que no tiende a $f'(0)=0$. Por el \cref{teo:heine}, $f'$ no es continua en~0.
  \item (c) $\abs{g(x)}\le\abs x\to0=g(0)$: $g$ es continua en~0. Pero $\frac{g(h)-g(0)}{h}=\sin\frac1h$, que no tiene límite en~0 (\cref{ej:sin1x}): $g$ no es derivable en~0. \hfill$\blacksquare$
\end{enumerate}
\resultado $f$ es derivable en $\R$ con $f'(0)=0$, pero $f'$ es discontinua en~0; $g$ es continua y no derivable en~0.
\comprobacion con $y_n=\frac{1}{2\pi n+\pi}$, $f'(y_n)=-\cos(2\pi n+\pi)=1$: $f'$ toma los valores $\pm1$ arbitrariamente cerca de~0, luego ni siquiera tiene límite allí. Aun así, $f'$ cumple la propiedad de los valores intermedios (\cref{eje:5-darboux}).
\end{solucion}

\begin{solucion}{eje:4-cadena-n}
\datos $F=f_n\circ\cdots\circ f_1$, $x_k=f_k(x_{k-1})$.
\planteamiento inducción sobre $n$ con el \cref{teo:cadena}; cota con la \cref{prop:derivadas-activaciones}.
\pasos
\begin{enumerate}[label=\arabic*.]
  \item Sea $F_n=f_n\circ\cdots\circ f_1$. Para $n=1$, $F_1'(x)=f_1'(x_0)$.
  \item Supóngase la fórmula para $F_n$. Como $F_{n+1}=f_{n+1}\circ F_n$ y $F_n(x)=x_n$, la regla de la cadena da $F_{n+1}'(x)=f_{n+1}'(x_n)\,F_n'(x)=f_{n+1}'(x_n)\prod_{k=1}^{n}f_k'(x_{k-1})=\prod_{k=1}^{n+1}f_k'(x_{k-1})$.
  \item Si todas las $f_k$ son $\sigma$, cada factor es $\sigma'(x_{k-1})\in(0,\frac14]$, luego $0<F'(x)\le(\frac14)^{n}=4^{-n}$. \hfill$\blacksquare$
\end{enumerate}
\resultado $F'(x)=\prod_kf_k'(x_{k-1})$; para $n$ sigmoides, $0<F'\le4^{-n}$.
\comprobacion e interpretación: con $n=10$, $F'\le4^{-10}\approx9{,}5\cdot10^{-7}$. En una red profunda, el gradiente respecto de los parámetros de las primeras capas contiene un producto así y se hace minúsculo: es el desvanecimiento del gradiente que motiva las activaciones ReLU y las conexiones residuales (módulo~11).
\end{solucion}
% ---------------------------------------------------------------------
\section{Tema 5}

\begin{solucion}{eje:5-hipotesis}
\datos cuatro funciones en intervalos.
\planteamiento comprobar continuidad en $[a,b]$ y derivabilidad en $(a,b)$; calcular la pendiente de la cuerda y resolver $f'(c)=m$.
\pasos
\begin{enumerate}[label=(\alph*)]
  \item No es derivable en $\frac12\in(-1,1)$: \textbf{no} se cumplen las hipótesis. Cuerda: $m=\frac{\frac12-\frac32}{2}=-\frac12$; pero $f'=\pm1$ donde existe. \textbf{No hay $c$}.
  \item Discontinua en $0\in[-1,1]$: \textbf{no} se cumplen. $m=\frac{1-1}{2}=0$ y $f'(x)=-2/x^{3}\ne0$. \textbf{No hay $c$}.
  \item \textbf{Sí} se cumplen. $m=e-1$ y $e^{c}=e-1$ da $c=\ln(e-1)\approx0{,}541\in(0,1)$.
  \item Continua, pero no derivable en $0$ (tangente vertical): \textbf{no} se cumplen. Sin embargo, $m=\frac{1-(-1)}{2}=1$ y $\frac13c^{-2/3}=1$ da $c=\pm3^{-3/2}\approx\pm0{,}192$: \textbf{sí hay} $c$.
\end{enumerate}
\resultado hipótesis solo en (c); hay $c$ en (c) y en (d), no en (a) ni en (b).
\comprobacion (d) muestra que las hipótesis son suficientes, no necesarias: el teorema no dice nada cuando fallan, y la conclusión puede cumplirse o no.
\end{solucion}

\begin{solucion}{eje:5-valores-c}
\datos cuatro funciones derivables en sus intervalos.
\planteamiento resolver $f'(c)=\frac{f(b)-f(a)}{b-a}$ y quedarse con las soluciones en $(a,b)$.
\pasos
\begin{enumerate}[label=(\alph*)]
  \item $m=1$, $2c=1$, $c=\frac12$.
  \item $m=\ln2$, $\frac1c=\ln2$, $c=\frac{1}{\ln2}\approx1{,}443\in(1,2)$.
  \item $m=\frac{6-(-6)}{4}=3$, $3c^{2}-1=3$, $c=\pm\frac{2}{\sqrt3}\approx\pm1{,}155$; ambos en $(-2,2)$.
  \item $m=\frac82=4$, $3c^{2}=4$, $c=\frac{2}{\sqrt3}\approx1{,}155$ (el valor negativo no está en $(0,2)$).
\end{enumerate}
\resultado (a) $\frac12$; (b) $\frac{1}{\ln2}$; (c) $\pm\frac{2}{\sqrt3}$; (d) $\frac{2}{\sqrt3}$.
\comprobacion (b): $\ln2\approx0{,}693$ y $1/1{,}443\approx0{,}693$.
\end{solucion}

\begin{solucion}{eje:5-lhopital}
\datos ocho límites indeterminados.
\planteamiento identificar la forma, reducirla a $0/0$ o $\infty/\infty$ y aplicar el \cref{teo:lhopital} (varias veces si hace falta).
\pasos
\begin{enumerate}[label=(\alph*)]
  \item $0/0$: $\frac{e^{x}-1}{2x}\to\frac12$.
  \item $0/0$: $\frac{\cos x-\cos^{-2}x}{3x^{2}}$, de nuevo $0/0$: $\frac{-\sin x-2\cos^{-3}x\sin x}{6x}=-\frac{\sin x}{x}\cdot\frac{1+2\cos^{-3}x}{6}\to-\frac{3}{6}=-\frac12$.
  \item $0/0$: $\frac{1}{1/x}=x\to1$.
  \item $\infty/\infty$: $\frac{x^{2}}{e^{x}}\to\frac{2x}{e^{x}}\to\frac{2}{e^{x}}\to0$.
  \item $0^{0}$: $x^{2x}=e^{2x\ln x}$ y $x\ln x\to0$ (\cref{ej:lhopital}(c)); el límite es $e^{0}=1$.
  \item $1^{\infty}$: $(1-\frac1x)^{x}=\exp\bigl(x\ln(1-\frac1x)\bigr)$. Con $t=1/x\to0^{+}$, $\frac{\ln(1-t)}{t}\overset{0/0}{\to}\frac{-1/(1-t)}{1}\to-1$. El límite es $e^{-1}$.
  \item $\infty-\infty$: $\frac{x-\sin x}{x\sin x}\overset{0/0}{\to}\frac{1-\cos x}{\sin x+x\cos x}\overset{0/0}{\to}\frac{\sin x}{2\cos x-x\sin x}\to\frac02=0$.
  \item $0/0$: $\frac{3^{x}\ln3-2^{x}\ln2}{1}\to\ln3-\ln2=\ln\frac32$.
\end{enumerate}
\resultado (a) $\frac12$; (b) $-\frac12$; (c) $1$; (d) $0$; (e) $1$; (f) $e^{-1}$; (g) $0$; (h) $\ln\frac32\approx0{,}405$.
\comprobacion (a), (b), (g) con desarrollos de Taylor (\cref{tema:6}): $e^{x}-x-1=\frac{x^{2}}{2}+o(x^{2})$; $\sin x-\tan x=-\frac{x^{3}}{2}+o(x^{3})$; $\frac1{\sin x}-\frac1x=\frac x6+o(x)$.
\end{solucion}

\begin{solucion}{eje:5-grafica}
\datos (a) $f(x)=\frac{x^{2}}{x-1}$; (b) $g(x)=x-3x^{1/3}$; (c) $h(x)=e^{-x^{2}}$.
\planteamiento dominio, asíntotas (\cref{def:asintotas}), signo de la primera derivada (crecimiento y extremos) y de la segunda (convexidad).
\pasos
\begin{enumerate}[label=(\alph*)]
  \item Dominio $\R\setminus\{1\}$. $f(x)=x+1+\frac{1}{x-1}$: asíntota vertical $x=1$ y oblicua $y=x+1$. $f'(x)=\frac{x(x-2)}{(x-1)^{2}}$: $f$ crece en $(-\infty,0]$ y en $[2,\infty)$ y decrece en $[0,1)$ y en $(1,2]$; máximo local $f(0)=0$ y mínimo local $f(2)=4$. $f''(x)=\frac{2}{(x-1)^{3}}$: cóncava en $(-\infty,1)$ y convexa en $(1,\infty)$; sin inflexiones (el cambio se produce en la asíntota).
  \item Dominio $\R$; $g$ es impar. $g'(x)=1-x^{-2/3}$ para $x\ne0$, negativa si $0<\abs x<1$ y positiva si $\abs x>1$: crece en $(-\infty,-1]$ y $[1,\infty)$, decrece en $[-1,1]$; máximo local $g(-1)=2$, mínimo local $g(1)=-2$. En $0$, $g'(x)\to-\infty$: tangente vertical. $g''(x)=\frac23x^{-5/3}$: cóncava para $x<0$, convexa para $x>0$; inflexión en $(0,0)$. Ceros: $x=0$ y $x^{2}=27$, $x=\pm3\sqrt3\approx\pm5{,}20$. Sin asíntotas: $g(x)/x\to1$ pero $g(x)-x=-3x^{1/3}$ no tiene límite finito.
  \item Dominio $\R$, par, $0<h\le1$. Asíntota horizontal $y=0$. $h'(x)=-2xe^{-x^{2}}$: crece en $(-\infty,0]$, decrece en $[0,\infty)$; máximo absoluto $h(0)=1$. $h''(x)=(4x^{2}-2)e^{-x^{2}}$: convexa si $\abs x>\frac{1}{\sqrt2}$ y cóncava si $\abs x<\frac{1}{\sqrt2}$; inflexiones en $\bigl(\pm\frac{1}{\sqrt2},e^{-1/2}\bigr)\approx(\pm0{,}707;\,0{,}607)$.
\end{enumerate}
\resultado los datos anteriores determinan las tres gráficas: (a) dos ramas, con un máximo local en el origen y un mínimo local en $(2,4)$; (b) una «S» tumbada con extremos en $(\mp1,\pm2)$ y tangente vertical en el origen; (c) la campana de Gauss.
\comprobacion (b): $g(-1)=-1-3(-1)=2$ y $g(1)=1-3=-2$, coherente con la imparidad. (c): $h$ es, salvo constantes, la densidad normal del módulo~5.
\end{solucion}

\begin{solucion}{eje:5-optimizacion}
\datos cuatro problemas de extremos.
\planteamiento reducir a una variable, puntos críticos, y justificar que son extremos absolutos (signo de la derivada o \cref{prop:intervalo-cerrado}).
\pasos
\begin{enumerate}[label=(\alph*)]
  \item $x=a-2y$, $P(y)=(a-2y)y=ay-2y^{2}$; $P'(y)=a-4y=0$ en $y=\frac a4$, y $P''=-4<0$; como $P$ es una parábola cóncava, es el máximo absoluto. $x=\frac a2$, $P=\frac{a^{2}}{8}$.
  \item $f(x)=-x^{2}\ln x$; $f'(x)=-2x\ln x-x=-x(2\ln x+1)$, positiva si $x<e^{-1/2}$ y negativa si $x>e^{-1/2}$: máximo absoluto en $x=e^{-1/2}$, con $f=-e^{-1}\cdot(-\frac12)=\frac{1}{2e}\approx0{,}184$.
  \item $k(x)=(x+1)e^{-x}$; $k'(x)=e^{-x}-(x+1)e^{-x}=-xe^{-x}<0$ para $x>0$: $k$ es estrictamente decreciente en $[0,\infty)$. Máximo absoluto $k(0)=1$. El ínfimo es $\lim_{x\to\infty}k(x)=0$ (\cref{prop:crecimientos}), pero $k(x)>0$ siempre: \textbf{no hay mínimo}.
  \item Con el diámetro sobre el eje $x$ y centro en el origen, el rectángulo tiene base $2x$ y altura $\sqrt{a^{2}-x^{2}}$, $x\in[0,a]$. Se maximiza $A^{2}=4x^{2}(a^{2}-x^{2})$, función de $u=x^{2}\in[0,a^{2}]$: $4u(a^{2}-u)$, máxima en $u=\frac{a^{2}}{2}$. Luego $x=\frac{a}{\sqrt2}$, base $a\sqrt2$, altura $\frac{a}{\sqrt2}$, área $a^{2}$; en los extremos $x=0$, $x=a$ el área es nula.
\end{enumerate}
\resultado (a) $x=\frac a2$, $y=\frac a4$, máximo $\frac{a^{2}}{8}$; (b) $\frac{1}{2e}$ en $x=e^{-1/2}$; (c) máximo $1$ en $x=0$, sin mínimo; (d) base $a\sqrt2$, altura $a/\sqrt2$, área $a^{2}$.
\comprobacion (d) con $a=1$: $x=0{,}7$ da $2\cdot0{,}7\cdot\sqrt{0{,}51}\approx0{,}9998<1$; $x=0{,}72$ da $0{,}9993<1$.
\end{solucion}

\begin{solucion}{eje:5-cubica}
\datos $y=x^{3}+ax^{2}+bx+c$, $y'=3x^{2}+2ax+b$.
\planteamiento la cúbica tiene un máximo y un mínimo locales si y solo si $y'$ cambia de signo dos veces, es decir, si tiene dos raíces reales distintas.
\pasos
\begin{enumerate}[label=\arabic*.]
  \item Si el discriminante $4a^{2}-12b>0$, $y'$ tiene raíces $x_1<x_2$ y su signo es $+,-,+$: por el criterio de la primera derivada, máximo local en $x_1$ y mínimo local en $x_2$.
  \item Si $4a^{2}-12b\le0$, $y'\ge0$ en todo $\R$ (con a lo sumo un cero), $y$ es estrictamente creciente (\cref{cor:tvm}(iii), aplicado a cada lado del cero) y no tiene extremos.
  \item La condición es, pues, $a^{2}>3b$; $c$ no interviene (traslada la gráfica en vertical).
  \item Ejemplos: $x^{3}-3x+1$ ($a=0$, $b=-3$): $0>-9$, extremos en $x=\pm1$. $x^{3}+x$ ($a=0$, $b=1$): $0>3$ es falso; $y'=3x^{2}+1>0$, sin extremos.
\end{enumerate}
\resultado tiene máximo y mínimo locales si y solo si $a^{2}>3b$.
\comprobacion criterio equivalente: el punto de inflexión está en $x_0=-\frac a3$, y la condición $a^{2}>3b$ equivale a $y'(x_0)=b-\frac{a^{2}}{3}<0$: la pendiente en la inflexión es negativa.
\end{solucion}

\begin{solucion}{eje:5-una-raiz}
\datos $p(x)=x^{3}+3x^{2}+16$; $h(x)=e^{x}+x-2$.
\planteamiento existencia con Bolzano; unicidad con la monotonía (\cref{cor:tvm}(iii)).
\pasos
\begin{enumerate}[label=\arabic*.]
  \item $p'(x)=3x(x+2)$: $p$ crece en $(-\infty,-2]$, decrece en $[-2,0]$ y crece en $[0,\infty)$. Por tanto, en $[-2,\infty)$ se tiene $p(x)\ge p(0)=16>0$, y no hay raíces en ese intervalo.
  \item En $(-\infty,-2]$, $p$ es estrictamente creciente, $p(-2)=20>0$ y $p(x)\to-\infty$: hay exactamente una raíz. Como $p(-4)=-64+48+16=0$, es $x=-4$.
  \item $h'(x)=e^{x}+1>0$: $h$ es estrictamente creciente y tiene a lo sumo un cero. $h(0)=-1<0<e-1=h(1)$: por Bolzano, exactamente uno, en $(0,1)$. \hfill$\blacksquare$
\end{enumerate}
\resultado $x^{3}+3x^{2}+16=0$ tiene la única raíz real $-4$; $e^{x}=2-x$ tiene una única solución, en $(0,1)$.
\comprobacion $p(x)=(x+4)(x^{2}-x+4)$, y $x^{2}-x+4$ tiene discriminante $1-16<0$. La solución de $e^{x}=2-x$ es $x\approx0{,}4429$.
\end{solucion}

\begin{solucion}{eje:5-desigualdades}
\datos cuatro desigualdades.
\planteamiento \cref{teo:tvm} en $[0,x]$ (o entre $x$ e $y$) y acotación de $f'(c)$.
\pasos
\begin{enumerate}[label=\arabic*.]
  \item (a) Si $x>1$, $\sin x\le1<x$. Si $0<x\le1$, el TVM da $\sin x-\sin0=x\cos c$ con $c\in(0,x)\subseteq(0,1)$; ahí $\cos c<1$, luego $\sin x<x$.
  \item (b) $\sqrt{1+x}-1=\frac{x}{2\sqrt{1+c}}$ con $c\in(0,x)$; como $\sqrt{1+c}>1$, $\sqrt{1+x}-1<\frac x2$.
  \item (c) $\ln(1+x)-\ln1=\frac{x}{1+c}$ con $c\in(0,x)$, y $\frac{1}{1+x}<\frac{1}{1+c}<1$; multiplicando por $x>0$, $\frac{x}{1+x}<\ln(1+x)<x$.
  \item (d) $\abs{\arctan't}=\frac{1}{1+t^{2}}\le1$, y el \cref{cor:tvm}(iv) da $\abs{\arctan x-\arctan y}\le\abs{x-y}$. \hfill$\blacksquare$
\end{enumerate}
\resultado las cuatro desigualdades.
\comprobacion con $x=0{,}5$: $\sin0{,}5=0{,}479<0{,}5$; $\sqrt{1{,}5}=1{,}2247<1{,}25$; $0{,}3333<\ln1{,}5=0{,}4055<0{,}5$.
\end{solucion}

\begin{solucion}{eje:5-comparacion}
\datos $f(a)\ge g(a)$, $f'>g'$ en $(a,\infty)$.
\planteamiento aplicar el \cref{cor:tvm}(iii) a $h=f-g$; después, inducción.
\pasos
\begin{enumerate}[label=\arabic*.]
  \item (a) $h=f-g$ es continua en $[a,\infty)$ y $h'>0$ en $(a,\infty)$, luego es estrictamente creciente en $[a,\infty)$. Para $x>a$, $h(x)>h(a)\ge0$, es decir, $f(x)>g(x)$.
  \item (b) Primero, $e^{x}>1=e^{0}$ para $x>0$, porque $\exp$ es estrictamente creciente. Con $f=e^{x}$ y $g=1+x$: $f(0)=g(0)=1$ y $f'=e^{x}>1=g'$ para $x>0$; por (a), $e^{x}>1+x$.
  \item Con $f=e^{x}$ y $g=1+x+\frac{x^{2}}{2}$: $f(0)=g(0)$ y $f'=e^{x}>1+x=g'$ para $x>0$ por el paso anterior; por (a), $e^{x}>1+x+\frac{x^{2}}{2}$.
  \item (c) Sea $P_n(x)=\sum_{k=0}^{n}\frac{x^{k}}{k!}$; nótese que $P_{n+1}'=P_n$ y $P_{n+1}(0)=1$. Si $e^{x}>P_n(x)$ para $x>0$, (a) con $f=e^{x}$, $g=P_{n+1}$ da $e^{x}>P_{n+1}(x)$ para $x>0$. El caso $n=0$ es el primer paso de (b). \hfill$\blacksquare$
\end{enumerate}
\resultado $e^{x}>\sum_{k=0}^{n}x^{k}/k!$ para todo $x>0$ y todo $n$.
\comprobacion con $x=1$, $n=3$: $e=2{,}718>1+1+0{,}5+0{,}1667=2{,}667$. El \cref{teo:taylor} da la misma conclusión: el resto es $\frac{e^{\xi}x^{n+1}}{(n+1)!}>0$.
\end{solucion}

\begin{solucion}{eje:5-raices-polinomio}
\datos $p$ de grado $n$ con raíces $x_1<\dots<x_n$; $f$ dos veces derivable con ceros $a<b<c$.
\planteamiento Rolle entre ceros consecutivos.
\pasos
\begin{enumerate}[label=\arabic*.]
  \item (a) Para $i=1,\dots,n-1$, $p(x_i)=p(x_{i+1})=0$ y Rolle da $c_i\in(x_i,x_{i+1})$ con $p'(c_i)=0$. Los intervalos son disjuntos, luego los $c_i$ son distintos: $p'$ tiene al menos $n-1$ raíces reales distintas, y no más, por ser de grado $n-1$.
  \item (b) Rolle en $[a,b]$ y en $[b,c]$ da $p_1\in(a,b)$ y $p_2\in(b,c)$ con $f'(p_1)=f'(p_2)=0$, y $p_1<p_2$. Rolle aplicado a $f'$ (derivable) en $[p_1,p_2]$ da $p\in(p_1,p_2)\subseteq(a,c)$ con $f''(p)=0$.
  \item (c) $q(x)=x^{5}+x-1$ tiene $q'(x)=5x^{4}+1>0$. Si tuviera dos raíces, Rolle daría un cero de $q'$; luego tiene a lo sumo una. Como $q(0)=-1<0<1=q(1)$, Bolzano da exactamente una, en $(0,1)$. \hfill$\blacksquare$
\end{enumerate}
\resultado (a) $p'$ tiene $n-1$ raíces reales distintas; (b) existe $p$; (c) una única raíz real.
\comprobacion (b) con $f=(x-a)(x-b)(x-c)$: $f''(x)=6x-2(a+b+c)$ se anula en la media $\frac{a+b+c}{3}\in(a,c)$. (c): la raíz es $x\approx0{,}7549$.
\end{solucion}

\begin{solucion}{eje:5-tvm-demo}
\datos $f$ dos veces derivable en $\R$.
\planteamiento (a) como en el \cref{teo:tvm}; (b) regla de L'Hôpital en la variable $h$, una sola vez, y la definición de $f''(a)$.
\pasos
\begin{enumerate}[label=\arabic*.]
  \item (a) \emph{Rolle:} si $f$ es continua en $[a,b]$, derivable en $(a,b)$ y $f(a)=f(b)$, existe $c\in(a,b)$ con $f'(c)=0$. \emph{TVM:} con $m=\frac{f(b)-f(a)}{b-a}$, la función $g(x)=f(x)-m(x-a)$ cumple las hipótesis de Rolle, pues $g(a)=g(b)=f(a)$; luego existe $c$ con $g'(c)=f'(c)-m=0$.
  \item (b) Sean $N(h)=f(a+h)+f(a-h)-2f(a)$ y $D(h)=h^{2}$. Ambas son derivables, $N(h)\to0$ y $D(h)\to0$ cuando $h\to0$, y $D'(h)=2h\ne0$ para $h\ne0$.
  \item $\frac{N'(h)}{D'(h)}=\frac{f'(a+h)-f'(a-h)}{2h}=\frac12\Bigl[\frac{f'(a+h)-f'(a)}{h}+\frac{f'(a-h)-f'(a)}{-h}\Bigr]\to\frac12\bigl[f''(a)+f''(a)\bigr]=f''(a)$, por la definición de $f''(a)$ como derivada de $f'$ (en el segundo cociente, con incremento $-h\to0$).
  \item Por el \cref{teo:lhopital}, $\frac{N(h)}{D(h)}\to f''(a)$. \hfill$\blacksquare$
\end{enumerate}
\resultado $f''(a)=\lim_{h\to0}\frac{f(a+h)+f(a-h)-2f(a)}{h^{2}}$.
\comprobacion con $f=\sin$, $a=1$, $h=0{,}01$: el cociente vale $-0{,}841464$, frente a $f''(1)=-\sin1=-0{,}841471$. Aplicar L'Hôpital dos veces exigiría que $f''$ fuera continua en $a$, hipótesis que no se tiene.
\end{solucion}

\begin{solucion}{eje:5-cauchy}
\datos $h(x)=f(x)-kg(x)$ con $k=\frac{f(b)-f(a)}{g(b)-g(a)}$.
\planteamiento (a) cálculo directo; (b) pendiente de la tangente a una curva paramétrica; (c) como en la demostración del \cref{teo:lhopital}.
\pasos
\begin{enumerate}[label=\arabic*.]
  \item (a) $h(b)-h(a)=\bigl(f(b)-f(a)\bigr)-k\bigl(g(b)-g(a)\bigr)=0$, por la definición de $k$.
  \item (b) La curva $t\mapsto(g(t),f(t))$, $t\in[a,b]$, une $(g(a),f(a))$ con $(g(b),f(b))$; la cuerda tiene pendiente $k$, y la tangente en $t=c$ tiene vector director $(g'(c),f'(c))$ y pendiente $f'(c)/g'(c)$. El teorema afirma que hay un punto de la curva con tangente paralela a la cuerda; con $g(t)=t$ es el TVM.
  \item (c) Sean $f(x),g(x)\to0$ cuando $x\to a^{+}$ y $\frac{f'}{g'}\to\ell$. Se define $f(a)=g(a)=0$. Para $x>a$ próximo, $g(x)\ne0$: si fuera $g(x)=0=g(a)$, Rolle daría un cero de $g'$ en $(a,x)$. Por el \cref{teo:cauchy} en $[a,x]$, $\frac{f(x)}{g(x)}=\frac{f'(c_x)}{g'(c_x)}$ con $a<c_x<x$, y cuando $x\to a^{+}$, $c_x\to a^{+}$ y el cociente tiende a $\ell$. \hfill$\blacksquare$
\end{enumerate}
\resultado $h(a)=h(b)$; interpretación geométrica; regla de L'Hôpital.
\comprobacion la hipótesis $g'\ne0$ se usa dos veces: para que $g(x)\ne0$ (y el cociente $f/g$ tenga sentido) y para que $f'(c_x)/g'(c_x)$ esté definido.
\end{solucion}

\begin{solucion}{eje:5-darboux}
\datos $g$ derivable cerca de $[a,b]$, $g'(a)<k<g'(b)$.
\planteamiento Weierstrass para $\varphi(x)=g(x)-kx$ y Fermat en el mínimo.
\pasos
\begin{enumerate}[label=\arabic*.]
  \item $\varphi$ es continua en $[a,b]$; por Weierstrass alcanza su mínimo en algún $c\in[a,b]$.
  \item $\varphi'(a)=g'(a)-k<0$. Como $\frac{\varphi(a+h)-\varphi(a)}{h}\to\varphi'(a)<0$, para $h>0$ pequeño $\varphi(a+h)<\varphi(a)$: el mínimo no está en $a$.
  \item $\varphi'(b)=g'(b)-k>0$. Como $\frac{\varphi(b-h)-\varphi(b)}{-h}\to\varphi'(b)>0$, para $h>0$ pequeño $\varphi(b-h)<\varphi(b)$: el mínimo no está en $b$.
  \item Luego $c\in(a,b)$ y, por el \cref{teo:fermat}, $\varphi'(c)=0$, es decir, $g'(c)=k$.
  \item \emph{Saltos.} Supóngase que $g'$ tiene en $x_0$ límites laterales $L_-<L_+$ (el caso $L_->L_+$ es análogo). Se eligen $k_1<k_2$ en $(L_-,L_+)$ tales que $g'(x_0)\notin[k_1,k_2]$, y $k\in(k_1,k_2)$. Existe $\delta>0$ con $g'<k_1$ en $(x_0-\delta,x_0)$ y $g'>k_2$ en $(x_0,x_0+\delta)$. Tomando $a$ y $b$ en esos intervalos, Darboux da $c\in(a,b)$ con $g'(c)=k$; pero $c=x_0$ es imposible ($g'(x_0)\ne k$), $c<x_0$ da $g'(c)<k$ y $c>x_0$ da $g'(c)>k$. Contradicción.
  \item La función escalón solo toma los valores $0$ y~$1$; si fuera $g'$, Darboux en $[-1,1]$ con $k=\frac12$ daría $g'(c)=\frac12$. \hfill$\blacksquare$
\end{enumerate}
\resultado las derivadas cumplen la propiedad de los valores intermedios, aunque no sean continuas, y no tienen discontinuidades de salto.
\comprobacion la derivada del \cref{eje:4-x2sen} es discontinua en~0, pero su discontinuidad es esencial (oscila entre $-1$ y $1$), no de salto, en coherencia con este resultado.
\end{solucion}

\begin{solucion}{eje:5-creciente-derivada}
\datos $f$ creciente y derivable en $a$.
\planteamiento signo de los cocientes incrementales y paso al límite.
\pasos
\begin{enumerate}[label=\arabic*.]
  \item (a) Si $h>0$, $f(a+h)\ge f(a)$ y el cociente $\frac{f(a+h)-f(a)}{h}$ es $\ge0$; si $h<0$, numerador $\le0$ y denominador $<0$, y el cociente también es $\ge0$. El límite de una función $\ge0$ es $\ge0$: $f'(a)\ge0$.
  \item (b) Si $f$ es estrictamente creciente, los cocientes son $>0$, pero el paso al límite solo conserva desigualdades no estrictas. Contraejemplo: $f(x)=x^{3}$ es estrictamente creciente y $f'(0)=0$.
  \item (c) $\sigma'=\sigma(1-\sigma)>0$ y $\tanh'=1-\tanh^{2}>0$ en todo $\R$, luego son estrictamente crecientes (\cref{cor:tvm}(iii)). Si $x<y$, $\ReLU(x)=\max(0,x)\le\max(0,y)=\ReLU(y)$, luego $\ReLU$ es creciente; no estrictamente, porque $\ReLU(-2)=\ReLU(-1)=0$. \hfill$\blacksquare$
\end{enumerate}
\resultado $f'(a)\ge0$; la desigualdad puede ser una igualdad aunque $f$ sea estrictamente creciente.
\comprobacion el recíproco parcial es el \cref{cor:tvm}(iii): $f'>0$ en un intervalo implica crecimiento estricto; $x^{3}$ muestra que la condición es suficiente pero no necesaria.
\end{solucion}

\begin{solucion}{eje:5-max-convexas}
\datos la \cref{def:convexa}.
\planteamiento encadenar desigualdades con $z=tx+(1-t)y$.
\pasos
\begin{enumerate}[label=\arabic*.]
  \item (a) Sea $h=\max(f,g)$ con $f,g$ convexas, $x,y\in I$, $t\in[0,1]$, $z=tx+(1-t)y$. Entonces $f(z)\le tf(x)+(1-t)f(y)\le th(x)+(1-t)h(y)$, y lo mismo para $g(z)$; luego $h(z)\le th(x)+(1-t)h(y)$.
  \item Las funciones afines $\ell(x)=mx+n$ son convexas (la desigualdad es una igualdad). Como $\ReLU(x)=\max(0,x)$ y $\abs x=\max(x,-x)$, ambas son convexas.
  \item (b) Sumando las desigualdades de $f$ y $g$ en $z$ se obtiene la de $f+g$.
  \item Si $f$ es convexa y $g$ convexa y creciente: $g(f(z))\le g\bigl(tf(x)+(1-t)f(y)\bigr)\le tg(f(x))+(1-t)g(f(y))$; la primera desigualdad usa que $g$ es creciente y $f$ convexa, y la segunda, que $g$ es convexa.
  \item (c) $x^{2}-1$ es convexa (derivada segunda $2>0$) y $\softplus$ es convexa ($\softplus''=\sigma'>0$) y creciente ($\softplus'=\sigma>0$). Por (b), $\softplus(x^{2}-1)$ es convexa. \hfill$\blacksquare$
\end{enumerate}
\resultado el máximo y la suma de convexas son convexas, y la composición con una convexa creciente también; $\ReLU$, $\abs x$ y $\softplus(x^{2}-1)$ son convexas.
\comprobacion la hipótesis de crecimiento es necesaria: $f(x)=x^{2}$ y $g(u)=-u$ (convexa, decreciente) dan $-x^{2}$, que no es convexa.
\end{solucion}

\begin{solucion}{eje:5-jensen}
\datos $f$ convexa en un intervalo $I$, $x_i\in I$, $t_i\ge0$, $\sum t_i=1$.
\planteamiento inducción sobre $n$, agrupando los $n$ primeros puntos.
\pasos
\begin{enumerate}[label=\arabic*.]
  \item Para $n=1$ es una igualdad, y para $n=2$ es la \cref{def:convexa}. Supóngase cierta para $n$ y sean $t_1,\dots,t_{n+1}$.
  \item Si $t_{n+1}=1$, todos los demás son nulos y la desigualdad es una igualdad. Si no, sea $s=1-t_{n+1}=\sum_{i\le n}t_i>0$ e $y=\sum_{i\le n}\frac{t_i}{s}x_i$, que es una combinación convexa de puntos de $I$ y, por tanto, está en $I$ (está entre el mínimo y el máximo de los $x_i$).
  \item Por la definición con dos puntos, $f\bigl(\sum_{i\le n+1}t_ix_i\bigr)=f(sy+t_{n+1}x_{n+1})\le sf(y)+t_{n+1}f(x_{n+1})$.
  \item Por la hipótesis de inducción con los pesos $t_i/s$, que suman~1, $f(y)\le\sum_{i\le n}\frac{t_i}{s}f(x_i)$. Sustituyendo, $f\bigl(\sum t_ix_i\bigr)\le\sum_{i\le n+1}t_if(x_i)$.
  \item (b) $-\ln$ es convexa en $(0,\infty)$, pues $(-\ln)''=1/x^{2}>0$. Con $t_i=\frac1n$: $-\ln\frac{\sum x_i}{n}\le-\frac1n\sum\ln x_i=-\ln\sqrt[n]{x_1\cdots x_n}$. Cambiando de signo y aplicando la exponencial, que es creciente, $\sqrt[n]{x_1\cdots x_n}\le\frac{x_1+\dots+x_n}{n}$. \hfill$\blacksquare$
\end{enumerate}
\resultado la desigualdad de Jensen y la desigualdad entre las medias aritmética y geométrica.
\comprobacion $x=(1,4,16)$: media geométrica $\sqrt[3]{64}=4$, aritmética $7$. Esta desigualdad, con $-\ln$, es la que en el módulo~1 relaciona la perplejidad con una media geométrica.
\end{solucion}
% ---------------------------------------------------------------------
\section{Tema 6}

\begin{solucion}{eje:6-coeficientes}
\datos $P_{2,0}f(x)=3-x+5x^{2}$; $p(x)=x^{3}$.
\planteamiento los coeficientes del polinomio de Taylor son $f^{(k)}(a)/k!$ (\cref{def:taylor}).
\pasos
\begin{enumerate}[label=\arabic*.]
  \item (a) $f(0)=3$; $f'(0)=-1$; $\frac{f''(0)}{2}=5$, luego $f''(0)=10$.
  \item Como $f'(0)=-1\ne0$, por el \cref{teo:fermat} $f$ no tiene extremo local en~0.
  \item (b) $p(1)=1$, $p'(1)=3$, $p''(1)=6$, $p'''(1)=6$: $P_{3,1}p(x)=1+3(x-1)+3(x-1)^{2}+(x-1)^{3}$.
  \item Por el binomio de Newton, $(1+(x-1))^{3}=1+3(x-1)+3(x-1)^{2}+(x-1)^{3}$, que es $x^{3}$.
\end{enumerate}
\resultado (a) $f(0)=3$, $f'(0)=-1$, $f''(0)=10$; no hay extremo en~0. (b) $1+3(x-1)+3(x-1)^{2}+(x-1)^{3}=x^{3}$.
\comprobacion en (b), con $x=2$: $1+3+3+1=8=2^{3}$.
\end{solucion}

\begin{solucion}{eje:6-por-que}
\datos la \cref{prop:taylor-unico} y el \cref{eje:4-par-impar}.
\planteamiento unicidad del polinomio de Taylor y paridad de las derivadas sucesivas.
\pasos
\begin{enumerate}[label=\arabic*.]
  \item Si $p$ es un polinomio de grado $\le n$, el propio $p$ es un polinomio de grado $\le n$ con las mismas derivadas que $p$ en $a$ hasta el orden $n$; por la unicidad, $P_{n,a}p=p$.
  \item Si $f$ es par, $f'$ es impar, $f''$ par, y en general $f^{(k)}$ es impar para $k$ impar. Una función impar definida en~0 se anula en~0, luego $f^{(k)}(0)=0$ para $k$ impar y solo quedan potencias pares.
  \item $\sin$ es impar, luego su polinomio de Maclaurin solo tiene potencias impares: el coeficiente de $x^{4}$ es $\sin^{(4)}(0)/4!=\sin0/24=0$, y $P_{4,0}\sin=P_{3,0}\sin$.
\end{enumerate}
\resultado las tres explicaciones anteriores.
\comprobacion $P_{4,0}\cos=1-\frac{x^{2}}{2}+\frac{x^{4}}{24}$ (solo potencias pares) y $P_{5,0}\tanh=x-\frac{x^{3}}{3}+\frac{2x^{5}}{15}$ (solo impares).
\end{solucion}

\begin{solucion}{eje:6-maclaurin}
\datos cinco funciones, $a=0$.
\planteamiento calcular $f^{(k)}(0)$ para $k\le n$ y dividir por $k!$.
\pasos
\begin{enumerate}[label=(\alph*)]
  \item $\cos$: derivadas en~0: $1,0,-1,0,1$. $P_4=1-\frac{x^{2}}{2}+\frac{x^{4}}{24}$.
  \item $e^{-x}$: $f^{(k)}(0)=(-1)^{k}$. $P_3=1-x+\frac{x^{2}}{2}-\frac{x^{3}}{6}$.
  \item $\frac{1}{1-x}$: $f^{(k)}(x)=\frac{k!}{(1-x)^{k+1}}$, $f^{(k)}(0)/k!=1$. $P_3=1+x+x^{2}+x^{3}$.
  \item $\sqrt{1+x}$: $f(0)=1$, $f'(0)=\frac12$, $f''(0)=-\frac14$. $P_2=1+\frac x2-\frac{x^{2}}{8}$.
  \item $\arctan$: $f'=\frac{1}{1+x^{2}}$, $f''=\frac{-2x}{(1+x^{2})^{2}}$, $f'''=\frac{6x^{2}-2}{(1+x^{2})^{3}}$; en~0: $0,1,0,-2$. $P_3=x-\frac{x^{3}}{3}$.
\end{enumerate}
\resultado los polinomios anteriores.
\comprobacion (b) es (a) del \cref{ej:maclaurin} con $x$ cambiado por $-x$; (d) es el binomio generalizado con $\alpha=\frac12$: $\binom{1/2}{2}=\frac{\frac12\cdot(-\frac12)}{2}=-\frac18$.
\end{solucion}

\begin{solucion}{eje:6-otros-puntos}
\datos $\ln x$ y $1/x$ en $a=1$; $\sqrt x$ en $a=4$.
\planteamiento derivadas en el punto y \cref{def:taylor}, en potencias de $(x-a)$.
\pasos
\begin{enumerate}[label=(\alph*)]
  \item $\ln$: en $x=1$, $f=0$, $f'=1$, $f''=-1$, $f'''=2$. $P_{3,1}=(x-1)-\frac{(x-1)^{2}}{2}+\frac{(x-1)^{3}}{3}$.
  \item $\sqrt x$: en $x=4$, $f=2$, $f'=\frac{1}{2\sqrt4}=\frac14$, $f''=-\frac{1}{4\cdot4^{3/2}}=-\frac{1}{32}$. $P_{2,4}=2+\frac{x-4}{4}-\frac{(x-4)^{2}}{64}$.
  \item $1/x$: $f^{(k)}(1)=(-1)^{k}k!$. $P_{3,1}=1-(x-1)+(x-1)^{2}-(x-1)^{3}$.
  \item $\ln1{,}1\approx0{,}1-0{,}005+0{,}000333=0{,}095333$, frente a $\ln1{,}1=0{,}095310$: error $2{,}3\cdot10^{-5}$.
\end{enumerate}
\resultado (a) $(x-1)-\frac{(x-1)^{2}}{2}+\frac{(x-1)^{3}}{3}$; (b) $2+\frac{x-4}{4}-\frac{(x-4)^{2}}{64}$; (c) $1-(x-1)+(x-1)^{2}-(x-1)^{3}$; $\ln1{,}1\approx0{,}09533$.
\comprobacion el resto de Lagrange de (a) es $\frac{f^{(4)}(\xi)}{24}(0{,}1)^{4}$ con $f^{(4)}(\xi)=-6/\xi^{4}$ y $\xi\in(1;1{,}1)$; su valor absoluto es menor que $\frac{6}{24}\cdot10^{-4}=2{,}5\cdot10^{-5}$, y el error real, $2{,}3\cdot10^{-5}$, lo respeta.
\end{solucion}

\begin{solucion}{eje:6-ln12}
\datos $f$ dos veces derivable; $f=\ln$, $a=1$, $x=1{,}2$.
\planteamiento $P_{2,a}$ y resto de orden~3 (\cref{teo:taylor} con $n=2$).
\pasos
\begin{enumerate}[label=\arabic*.]
  \item (a) $f(x)\approx f(a)+f'(a)(x-a)+\frac{f''(a)}{2}(x-a)^{2}$.
  \item (b) Con $\ln$ en $a=1$: $\ln(1+h)\approx h-\frac{h^{2}}{2}$; para $h=0{,}2$, $\ln1{,}2\approx0{,}2-0{,}02=0{,}18$.
  \item (c) $R_2=\frac{f'''(\xi)}{6}h^{3}$ con $f'''(\xi)=\frac{2}{\xi^{3}}$, $\xi\in(1;1{,}2)$. Como $\xi>1$, $0<R_2<\frac{2}{6}(0{,}2)^{3}=0{,}00267$.
  \item El resto es positivo: $0{,}18$ aproxima por defecto, y $0{,}18<\ln1{,}2<0{,}18267$.
\end{enumerate}
\resultado $\ln1{,}2\approx0{,}18$, con error menor que $0{,}0027$ y por defecto.
\comprobacion $\ln1{,}2=0{,}18232$; el error real es $0{,}00232$, dentro de la cota y del signo previstos.
\end{solucion}

\begin{solucion}{eje:6-cuadraticas}
\datos seis funciones cerca de~0.
\planteamiento $P_{2,0}$, o composición de desarrollos conocidos (\cref{tab:maclaurin}) truncando en grado~2.
\pasos
\begin{enumerate}[label=(\alph*)]
  \item $\tan0=0$, $\tan'0=1$, $\tan''x=2\tan x(1+\tan^{2}x)$ vale~0 en~0: $\tan x\approx x$ (el término cuadrático es nulo, por ser $\tan$ impar).
  \item $\ln\cos x=\ln\bigl(1-\frac{x^{2}}{2}+o(x^{2})\bigr)=-\frac{x^{2}}{2}+o(x^{2})$, usando $\ln(1+u)=u+o(u)$.
  \item $\sin2x\approx2x$; $\cos2x\approx1-\frac{(2x)^{2}}{2}=1-2x^{2}$.
  \item $\frac{1}{\cos x}=\frac{1}{1-x^{2}/2+o(x^{2})}=1+\frac{x^{2}}{2}+o(x^{2})$, usando $\frac{1}{1-u}=1+u+o(u)$.
  \item $e^{x^{2}}=1+x^{2}+o(x^{2})$.
  \item $\frac{1}{a+bx}=\frac1a\cdot\frac{1}{1+\frac bax}=\frac1a\Bigl(1-\frac bax+\frac{b^{2}}{a^{2}}x^{2}\Bigr)+o(x^{2})=\frac1a-\frac{b}{a^{2}}x+\frac{b^{2}}{a^{3}}x^{2}+o(x^{2})$.
\end{enumerate}
\resultado (a) $x$; (b) $-\frac{x^{2}}{2}$; (c) $2x$ y $1-2x^{2}$; (d) $1+\frac{x^{2}}{2}$; (e) $1+x^{2}$; (f) $\frac1a-\frac{b}{a^{2}}x+\frac{b^{2}}{a^{3}}x^{2}$.
\comprobacion (b) derivando: $(\ln\cos x)'=-\tan x$, $(\ln\cos x)''=-(1+\tan^{2}x)$, que vale $-1$ en~0: coeficiente $-\frac12$.
\end{solucion}

\begin{solucion}{eje:6-error}
\datos (a) $x=0{,}1$; (b) $x=0{,}5$, error $<10^{-5}$; (c) $\sqrt[3]{8{,}1}$, $a=8$.
\planteamiento \cref{cor:cota} con una cota de la derivada siguiente en el intervalo.
\pasos
\begin{enumerate}[label=(\alph*)]
  \item $\sin0{,}1\approx0{,}1-\frac{0{,}001}{6}=0{,}0998333$. Como $P_{3,0}\sin=P_{4,0}\sin$, se puede usar el resto de orden~5: $\abs R\le\frac{(0{,}1)^{5}}{120}\approx8{,}3\cdot10^{-8}$.
  \item $\abs{R_n}=\frac{e^{\xi}(0{,}5)^{n+1}}{(n+1)!}<\frac{2\cdot(0{,}5)^{n+1}}{(n+1)!}$, pues $e^{\xi}<e^{0{,}5}<2$. Para $n=5$ la cota es $4{,}3\cdot10^{-5}$; para $n=6$, $3{,}1\cdot10^{-6}<10^{-5}$. Basta $n=6$.
  \item $f(x)=x^{1/3}$, $f(8)=2$, $f'(8)=\frac13\cdot8^{-2/3}=\frac{1}{12}$: $\sqrt[3]{8{,}1}\approx2+\frac{0{,}1}{12}=2{,}008333$. $f''(t)=-\frac29t^{-5/3}$, y para $t\ge8$, $\abs{f''(t)}\le\frac29\cdot\frac{1}{32}=\frac{1}{144}$. Luego $\abs R\le\frac{1}{2}\cdot\frac{1}{144}\cdot(0{,}1)^{2}\approx3{,}47\cdot10^{-5}$.
\end{enumerate}
\resultado (a) $0{,}0998333$ con error $<8{,}4\cdot10^{-8}$; (b) $n=6$; (c) $2{,}008333$ con error $<3{,}5\cdot10^{-5}$ (por exceso, porque $f''<0$).
\comprobacion valores reales: $\sin0{,}1=0{,}09983342$ (error $8{,}3\cdot10^{-8}$); $P_{6,0}(0{,}5)=1{,}6487196$ frente a $e^{0{,}5}=1{,}6487213$ (error $1{,}7\cdot10^{-6}$); $\sqrt[3]{8{,}1}=2{,}0082989$ (error $3{,}45\cdot10^{-5}$).
\end{solucion}

\begin{solucion}{eje:6-limites}
\datos cinco límites en~0.
\planteamiento sustituir desarrollos hasta el primer término no nulo del numerador (\cref{tab:maclaurin}, con un término más cuando haga falta).
\pasos
\begin{enumerate}[label=(\alph*)]
  \item $e^{x}-1-x-\frac{x^{2}}{2}=\frac{x^{3}}{6}+o(x^{3})$: límite $\frac16$.
  \item $\cos x-1+\frac{x^{2}}{2}=\frac{x^{4}}{24}+o(x^{5})$: límite $\frac{1}{24}$.
  \item $\sin x-x\cos x=\bigl(x-\frac{x^{3}}{6}\bigr)-x\bigl(1-\frac{x^{2}}{2}\bigr)+o(x^{4})=\frac{x^{3}}{3}+o(x^{4})$: límite $\frac13$.
  \item $\tanh x-x=-\frac{x^{3}}{3}+o(x^{4})$: límite $-\frac13$.
  \item $\sigma(x)-\frac12-\frac x4=-\frac{x^{3}}{48}+o(x^{4})$ (\cref{ej:taylor-activaciones}): límite $-\frac{1}{48}$.
\end{enumerate}
\resultado (a) $\frac16$; (b) $\frac{1}{24}$; (c) $\frac13$; (d) $-\frac13$; (e) $-\frac{1}{48}$.
\comprobacion (c) con L'Hôpital: $\frac{\cos x-\cos x+x\sin x}{3x^{2}}=\frac{\sin x}{3x}\to\frac13$. (d) y (e) son coherentes entre sí: $\sigma(x)-\frac12=\frac12\tanh\frac x2$, y $\frac12\cdot\bigl(-\frac13\bigr)\cdot\frac18=-\frac{1}{48}$.
\end{solucion}

\begin{solucion}{eje:6-poligono}
\datos polígono regular de $n$ lados inscrito en la circunferencia unidad.
\planteamiento teorema del coseno en el triángulo formado por el centro y dos vértices consecutivos; aproximación $\cos\theta=1-\frac{\theta^{2}}{2}+o(\theta^{3})$.
\pasos
\begin{enumerate}[label=\arabic*.]
  \item (a) Dos radios de longitud~1 forman un ángulo central $\theta=\frac{2\pi}{n}$; el lado opuesto mide $\ell$ con $\ell^{2}=1+1-2\cos\theta$, es decir, $\ell=\sqrt{2-2\cos\frac{2\pi}{n}}$.
  \item (b) $2-2\cos\theta=\theta^{2}+o(\theta^{3})$, luego $\ell=\theta\sqrt{1+o(\theta)}$ y el perímetro es $n\ell=n\theta\sqrt{1+o(\theta)}=2\pi\sqrt{1+o(1/n)}\to2\pi$.
  \item (c) $n=6$: $\ell=\sqrt{2-2\cdot\frac12}=1$, perímetro $6$. $n=100$: perímetro $\approx6{,}28215$.
\end{enumerate}
\resultado el perímetro tiende a $2\pi$, la longitud de la circunferencia.
\comprobacion $2\pi\approx6{,}28319$; para $n=100$ la diferencia es $0{,}00103$, del orden de $\frac{2\pi\cdot\theta^{2}}{24}=\frac{(2\pi)^{3}}{24n^{2}}\approx0{,}00103$, que es lo que predice el término siguiente del desarrollo.
\end{solucion}

\begin{solucion}{eje:6-taylor-orden1}
\datos $F(t)=f(t)-f(a)-f'(a)(t-a)$, $G(t)=(t-a)^{2}$.
\planteamiento dos aplicaciones del \cref{teo:cauchy}.
\pasos
\begin{enumerate}[label=\arabic*.]
  \item Sea $x\ne a$ (si $x=a$ es trivial). $F(a)=G(a)=0$, y $G'(t)=2(t-a)\ne0$ para $t$ estrictamente entre $a$ y $x$.
  \item Por Cauchy entre $a$ y $x$: $\frac{F(x)}{G(x)}=\frac{F(x)-F(a)}{G(x)-G(a)}=\frac{F'(c_1)}{G'(c_1)}=\frac{f'(c_1)-f'(a)}{2(c_1-a)}$, con $c_1$ entre $a$ y $x$.
  \item Sean $F_1(t)=f'(t)-f'(a)$ y $G_1(t)=2(t-a)$, derivables con $G_1'=2\ne0$ y $F_1(a)=G_1(a)=0$. Por Cauchy entre $a$ y $c_1$: $\frac{F_1(c_1)}{G_1(c_1)}=\frac{f''(\xi)}{2}$, con $\xi$ entre $a$ y $c_1$, y por tanto entre $a$ y $x$.
  \item Luego $F(x)=\frac{f''(\xi)}{2}G(x)$, es decir, $f(x)=f(a)+f'(a)(x-a)+\frac12f''(\xi)(x-a)^{2}$.
  \item Si $f''\ge0$, el último término es $\ge0$ y $f(x)\ge f(a)+f'(a)(x-a)$: la gráfica está por encima de la tangente en $a$. \hfill$\blacksquare$
\end{enumerate}
\resultado el teorema de Taylor de orden~1 con resto de Lagrange, y la desigualdad de la tangente para $f''\ge0$.
\comprobacion con $f(x)=x^{3}$, $a=0$: $x^{3}=0+0+\frac12\cdot6\xi\cdot x^{2}$ da $\xi=\frac x3$, que está entre $0$ y $x$. La conclusión final reproduce el \cref{teo:convexidad}(i) por otra vía.
\end{solucion}

\begin{solucion}{eje:6-binomica}
\datos $f(x)=(1+x)^{1/2}$, $c_n=\binom{1/2}{n}$.
\planteamiento (b) cociente de coeficientes consecutivos e inducción; (c) acotar el resto de Lagrange.
\pasos
\begin{enumerate}[label=\arabic*.]
  \item (a) Si $f$ es $n+1$ veces derivable en un intervalo abierto que contiene a $a$ y $x$, existe $\xi$ entre $a$ y $x$ con $f(x)=\sum_{k=0}^{n}\frac{f^{(k)}(a)}{k!}(x-a)^{k}+\frac{f^{(n+1)}(\xi)}{(n+1)!}(x-a)^{n+1}$ (\cref{teo:taylor}).
  \item (b) $c_0=1$ y $c_{n+1}=c_n\cdot\frac{\frac12-n}{n+1}$. Como $\abs{\frac12-n}\le n+\frac12<n+1$, $\abs{c_{n+1}}\le\abs{c_n}$; por inducción, $\abs{c_n}\le1$.
  \item (c) $f^{(k)}(t)=\frac12\bigl(\frac12-1\bigr)\cdots\bigl(\frac12-k+1\bigr)(1+t)^{\frac12-k}=k!\,c_k(1+t)^{\frac12-k}$. Luego $f^{(k)}(0)/k!=c_k$ y, para $x\in(0,1)$, el resto es $R_n(x)=c_{n+1}(1+\xi)^{\frac12-n-1}x^{n+1}$ con $0<\xi<x$.
  \item Como $\frac12-n-1<0$ y $1+\xi>1$, $(1+\xi)^{\frac12-n-1}<1$; por (b), $\abs{R_n(x)}\le x^{n+1}$, que tiende a~0 porque $0<x<1$.
  \item Por tanto, $\sum_{k=0}^{n}c_kx^{k}=f(x)-R_n(x)\to\sqrt{1+x}$. \hfill$\blacksquare$
\end{enumerate}
\resultado $\sqrt{1+x}=\sum_{n\ge0}\binom{1/2}{n}x^{n}$ para $x\in(0,1)$.
\comprobacion los primeros coeficientes son $1,\frac12,-\frac18,\frac1{16}$; con $x=0{,}21$, $1+0{,}105-0{,}0055125+0{,}0005788=1{,}1000663$, y $\sqrt{1{,}21}=1{,}1$, con error del orden de $\abs{c_4}x^{4}\approx7{,}6\cdot10^{-5}$.
\end{solucion}

\begin{solucion}{eje:6-ln2}
\datos $f(x)=\ln(1+x)$, $x\in[0,1]$.
\planteamiento resto de Lagrange con $f^{(n+1)}(t)=\frac{(-1)^{n}n!}{(1+t)^{n+1}}$ (\cref{ej:maclaurin}(c)).
\pasos
\begin{enumerate}[label=\arabic*.]
  \item $R_n(x)=\frac{f^{(n+1)}(\xi)}{(n+1)!}x^{n+1}=\frac{(-1)^{n}}{(n+1)(1+\xi)^{n+1}}x^{n+1}$ con $0\le\xi\le x$.
  \item Como $(1+\xi)^{n+1}\ge1$, $\abs{R_n(x)}\le\frac{x^{n+1}}{n+1}$.
  \item Con $x=1$: $\bigl|\ln2-\sum_{k=1}^{n}\frac{(-1)^{k-1}}{k}\bigr|\le\frac{1}{n+1}\to0$, luego la suma parcial tiende a $\ln2$.
  \item Error menor que $10^{-3}$: basta $\frac{1}{n+1}<10^{-3}$, es decir, $n\ge1000$. \hfill$\blacksquare$
\end{enumerate}
\resultado $\ln2=1-\frac12+\frac13-\cdots$; la cota garantiza error $<10^{-3}$ con $1000$ términos.
\comprobacion la convergencia es muy lenta: con $n=10$, la suma es $0{,}6456$ y $\ln2=0{,}6931$, error $0{,}048<\frac1{11}$. Para calcular $\ln2$ en la práctica se usan otras series, como la de $\artanh$: $\ln2=2\artanh\frac13$.
\end{solucion}

\begin{solucion}{eje:6-e-irracional}
\datos $e=\sum_{k=0}^{n}\frac1{k!}+\frac{e^{\xi}}{(n+1)!}$ con $0<\xi<1$ (\cref{ej:calculo-e}); $2<e<3$.
\planteamiento reducción al absurdo: multiplicar por $q!$ y obtener un entero entre $0$ y $1$.
\pasos
\begin{enumerate}[label=\arabic*.]
  \item Supóngase $e=p/q$ con $p,q\in\N$; multiplicando numerador y denominador por~2 si hace falta, puede suponerse $q\ge2$.
  \item Con $n=q$: $q!\,e-\sum_{k=0}^{q}\frac{q!}{k!}=\frac{e^{\xi}}{q+1}$.
  \item El miembro izquierdo es entero: $q!\,e=(q-1)!\,p$, y $q!/k!$ es entero para $k\le q$.
  \item El derecho cumple $0<\frac{e^{\xi}}{q+1}<\frac{3}{q+1}\le1$, porque $q\ge2$.
  \item No hay enteros estrictamente entre $0$ y $1$: contradicción. Luego $e$ es irracional. \hfill$\blacksquare$
\end{enumerate}
\resultado $e\notin\Q$.
\comprobacion la cota $e<3$ se demostró en el \cref{teo:e}; la cota $e^{\xi}>0$ es la que impide que el resto sea nulo.
\end{solucion}

\begin{solucion}{eje:6-criterio-n}
\datos cuatro funciones con punto crítico en~0; dos desigualdades.
\planteamiento (a) primera derivada no nula en~0 y \cref{prop:criterio-n}; (b) resto de Lagrange de orden adecuado con signo conocido.
\pasos
\begin{enumerate}[label=\arabic*.]
  \item $x^{4}$: $f^{(4)}(0)=24>0$, $n=4$ par: mínimo estricto. $x^{5}$: $f^{(5)}(0)=120$, $n=5$ impar: no hay extremo.
  \item $f(x)=\cos x-1+\frac{x^{2}}{2}$: $f'=-\sin x+x$, $f''=1-\cos x$, $f'''=\sin x$ se anulan en~0, y $f^{(4)}(0)=\cos0=1>0$: mínimo estricto.
  \item $f(x)=x\sin x$: $f'=\sin x+x\cos x$ se anula en~0 y $f''=2\cos x-x\sin x$ vale $2>0$: mínimo estricto.
  \item (b) Para $x\ge0$, Taylor de orden~2 de $\sin$ en~0: $\sin x=x-\frac{\cos\xi}{6}x^{3}$, y como $\cos\xi\le1$ y $x^{3}\ge0$, $\sin x\ge x-\frac{x^{3}}{6}$.
  \item Para todo $x$, Taylor de orden~1 de $\cos$ en~0: $\cos x=1-\frac{\cos\xi}{2}x^{2}\ge1-\frac{x^{2}}{2}$. \hfill$\blacksquare$
\end{enumerate}
\resultado mínimos estrictos en~0 para $x^{4}$, $\cos x-1+\frac{x^{2}}{2}$ y $x\sin x$; ningún extremo para $x^{5}$; las dos desigualdades.
\comprobacion la desigualdad $\cos x\ge1-\frac{x^{2}}{2}$ dice que $\cos x-1+\frac{x^{2}}{2}\ge0=f(0)$: el mínimo del paso~2 es, de hecho, absoluto.
\end{solucion}

\begin{solucion}{eje:6-diferencias}
\datos $f\in C^{3}$, $\abs{f''}\le M_2$, $\abs{f'''}\le M_3$.
\planteamiento \cref{teo:taylor} de orden~1 para la diferencia progresiva y de orden~2 en $a+h$ y $a-h$ para la centrada.
\pasos
\begin{enumerate}[label=\arabic*.]
  \item $f(a+h)=f(a)+f'(a)h+\frac{f''(\xi)}{2}h^{2}$, luego $\frac{f(a+h)-f(a)}{h}-f'(a)=\frac{f''(\xi)}{2}h$ y su valor absoluto es $\le\frac{M_2}{2}\abs h$.
  \item $f(a\pm h)=f(a)\pm f'(a)h+\frac{f''(a)}{2}h^{2}\pm\frac{f'''(\xi_\pm)}{6}h^{3}$. Restando, los términos de orden~0 y~2 se cancelan: $f(a+h)-f(a-h)=2f'(a)h+\frac{f'''(\xi_+)+f'''(\xi_-)}{6}h^{3}$.
  \item Dividiendo por $2h$: $\frac{f(a+h)-f(a-h)}{2h}-f'(a)=\frac{f'''(\xi_+)+f'''(\xi_-)}{12}h^{2}$, cuyo valor absoluto es $\le\frac{2M_3}{12}h^{2}=\frac{M_3}{6}h^{2}$.
  \item Numéricamente, con $f=\sin$, $a=1$, $h=0{,}01$ y $M_2=M_3=1$: el error de la progresiva es $-4{,}22\cdot10^{-3}$ (cota $5\cdot10^{-3}$) y el de la centrada, $-9{,}0\cdot10^{-6}$ (cota $1{,}67\cdot10^{-5}$). \hfill$\blacksquare$
\end{enumerate}
\resultado errores $O(h)$ y $O(h^{2})$; la diferencia centrada es preferible: con el mismo $h$, es unas 500 veces más precisa en el ejemplo.
\comprobacion al dividir $h$ por~10, el error de la progresiva se divide por unos 10 y el de la centrada por unos 100; es el análisis de la comprobación de gradientes del módulo~4, donde además interviene el error de redondeo, que crece cuando $h$ disminuye.
\end{solucion}
